\documentclass[10pt,a4paper,oneside]{article}

% ---------------------------------------------------------------------------
% ml-setor-energetico-en.tex -- en-US version of `ml-setor-energetico.tex`.
%
% PARITY WITH THE PORTUGUESE FILE. This is a translation, not a fork: both
% documents carry the same sections, the same equations, the same numbers and
% the same references. A change in one belongs in the other.
%
% DELIBERATE DIVERGENCES FROM THE SOURCE FILE:
%
%   `babel` loads `american`. The cover block, the footer and the NOTE box
%   come from the shared preamble with Portuguese labels; they are redefined
%   below, locally, so the preamble keeps serving the Portuguese documents
%   unchanged.
%
%   Numbers follow en-US: decimal point, comma as thousands separator.
%
%   Agency names, program names, acronyms of Brazilian indicators (MMGD, PLD,
%   DEC, FEC, DIC, FIC, DMIC) and law numbers stay as in the source: they are
%   identifiers. Each one is glossed in English at first use. The superscript
%   PNT of the loss balance is NTL (non-technical losses).
%
% Base preamble of the repository: font, type area, packages, caption, title
% and header. Compile from inside `docs/product/`.
% ---------------------------------------------------------------------------
\usepackage[american]{babel}
% =====================================================================
%  Preâmbulo base: pacotes, tipografia e convenções de todo documento
%  LaTeX do repositório
%
%  Origem: `config/preambulo.tex` do documento de projeto do rhae-sensor.
%  As seções abaixo seguem a ordem de lá, e o que difere está marcado
%  com DIFERENÇA.
%
%  Convenções tipográficas e de apresentação adotadas
%    ISO 216        formato A4
%    ISO 80000-1    números e unidades (vírgula decimal, espaço fino entre
%                   grupos de três algarismos, unidade em tipo romano)
%    Tabelas sem filetes verticais (booktabs); legenda de tabela acima,
%    de figura abaixo.
%
%  Uso, antes do `\begin{document}`:
%    report/*.tex          via `\usepackage{estilo-reporte}`, que carrega este
%    docs/product/*.tex    % =====================================================================
%  Preâmbulo base: pacotes, tipografia e convenções de todo documento
%  LaTeX do repositório
%
%  Origem: `config/preambulo.tex` do documento de projeto do rhae-sensor.
%  As seções abaixo seguem a ordem de lá, e o que difere está marcado
%  com DIFERENÇA.
%
%  Convenções tipográficas e de apresentação adotadas
%    ISO 216        formato A4
%    ISO 80000-1    números e unidades (vírgula decimal, espaço fino entre
%                   grupos de três algarismos, unidade em tipo romano)
%    Tabelas sem filetes verticais (booktabs); legenda de tabela acima,
%    de figura abaixo.
%
%  Uso, antes do `\begin{document}`:
%    report/*.tex          via `\usepackage{estilo-reporte}`, que carrega este
%    docs/product/*.tex    % =====================================================================
%  Preâmbulo base: pacotes, tipografia e convenções de todo documento
%  LaTeX do repositório
%
%  Origem: `config/preambulo.tex` do documento de projeto do rhae-sensor.
%  As seções abaixo seguem a ordem de lá, e o que difere está marcado
%  com DIFERENÇA.
%
%  Convenções tipográficas e de apresentação adotadas
%    ISO 216        formato A4
%    ISO 80000-1    números e unidades (vírgula decimal, espaço fino entre
%                   grupos de três algarismos, unidade em tipo romano)
%    Tabelas sem filetes verticais (booktabs); legenda de tabela acima,
%    de figura abaixo.
%
%  Uso, antes do `\begin{document}`:
%    report/*.tex          via `\usepackage{estilo-reporte}`, que carrega este
%    docs/product/*.tex    \input{../../report/config/preambulo}
%
%  FORMA DE MANUAL TÉCNICO. Fonte, números, legendas e pacotes vêm da
%  origem. A forma da página não: é de manual de referência, e não de
%  artigo. Coluna única, títulos recuados para a esquerda do texto,
%  rodapé com identificação e "página X de Y" em toda página, avisos com
%  palavra de sinal (ANSI Z535.6), passos numerados e histórico de
%  revisões (IEC/IEEE 82079-1). O desenho de página segue a classe
%  `refman` do CTAN.
%
%  O documento carrega `babel` ANTES deste arquivo. Documento deitado, que
%  existe para caber tabela larga, carrega também `geometry` antes; aí a
%  margem de manual não entra e os títulos não avançam.
%
%  Motor: pdflatex. Charis SIL, newtxmath e Bera Mono entram como Type 1
%  em T1, e a matemática do newtxmath não passa pelo `unicode-math`.
% =====================================================================
\RequirePackage{iftex}
\RequirePDFTeX
% Este arquivo entra tanto de um `.tex` quanto de dentro do
% `estilo-reporte.sty`, onde o `@` ja e letra. O codigo de categoria e
% devolvido como estava, e nao forcado a "outro" no fim.
\edef\pbrestauraarroba{\catcode`\noexpand\@=\the\catcode`\@\relax}
\makeatletter

% ---------------------------------------------------------------------
%  Fontes, idioma e microtipografia
% ---------------------------------------------------------------------
\usepackage[T1]{fontenc}
\usepackage[utf8]{inputenc}
% DIFERENÇA: o idioma é do documento, porque há documento em inglês.
\@ifpackageloaded{babel}{}{\usepackage[brazilian]{babel}}
\usepackage{CharisSIL}
\usepackage[charter]{newtxmath}
\usepackage[scaled=0.83]{beramono}
\usepackage[protrusion=true,expansion=true,tracking=true]{microtype}
\SetTracking{encoding=*,shape=sc}{40}
% DIFERENÇA: a origem não usa sem serifa. Os documentos daqui pedem
% `\sffamily` em título, tabela e nota; sem esta linha ela cairia na
% Computer Modern Sans, que não pertence ao conjunto. Uma família só.
\renewcommand{\sfdefault}{\rmdefault}

% DIFERENÇA: mancha de manual, e não a coluna de 468 pt da origem. A4 com
% 22 mm de margem de cada lado, 166 mm úteis. O texto corre em 152 mm e
% os 14 mm da esquerda são o recuo dos títulos: o bastante para o título
% sair do alinhamento do texto, sem deixar coluna vazia ao lado dele. A
% primeira versão reservava 40 mm, e a margem ficava em branco em toda
% página que não fosse título; o documento crescia um terço só por isso.
% `\margemmanual` é quanto os títulos, o cabeçalho e o rodapé avançam
% para a esquerda; vale zero quando o documento traz mancha própria.
% A altura é de 253 mm, a mesma da mancha anterior, porque as tabelas
% deitadas foram dimensionadas para ela.
\newlength{\margemmanual}
\@ifpackageloaded{geometry}{\setlength{\margemmanual}{0pt}}{%
  \setlength{\margemmanual}{14mm}%
  \usepackage[a4paper,left=36mm,right=22mm,top=20mm,bottom=24mm,
              headheight=14pt,headsep=7mm,footskip=12mm]{geometry}}

\usepackage{amsmath}
\usepackage[dvipsnames,table]{xcolor}
\usepackage{graphicx}
\usepackage{tikz}
\usepackage{booktabs}
\usepackage{tabularx}
\usepackage{array}
\usepackage{xltabular}
\usepackage{float}
\usepackage{pdflscape}
\usepackage{rotating}
% Página com o conteúdo girado, mas exibida em retrato no leitor de PDF
\newenvironment{paisagemretrato}{\begingroup\def\PLS@AddRotate##1{}\landscape}{\endlandscape\endgroup}
\usepackage{enumitem}
\usepackage{changepage}
\usepackage[most]{tcolorbox}
\usepackage{caption}
\usepackage{titlesec}
\usepackage{fancyhdr}
\usepackage{lastpage}
\usepackage{csquotes}
\usepackage{siunitx}
% DIFERENÇA: sem `biblatex`. Nenhum documento daqui tem `.bib`; as
% referências são listas escritas à mão.
\usepackage{xurl}
\usepackage[hidelinks,pdfusetitle]{hyperref}

% ---------------------------------------------------------------------
%  Números e unidades (ISO 80000-1)
% ---------------------------------------------------------------------
\sisetup{
  output-decimal-marker = {,},
  group-separator       = {\,},
  group-minimum-digits  = 4,
  exponent-product      = \cdot,
  uncertainty-mode      = separate,
  range-phrase          = {\,a\,},
  list-pair-separator   = {\ e\ },
  list-final-separator  = {\ e\ },
  per-mode              = symbol,
}

% ---------------------------------------------------------------------
%  Parágrafo, listas e flutuantes
% ---------------------------------------------------------------------
% DIFERENÇA: parágrafo em bloco, separado por espaço e sem recuo, que é
% o de manual; a origem recua 1 em.
\setlength{\parindent}{0pt}
\setlength{\parskip}{0.45em plus 0.1em}
\linespread{1.04}
\setlist{itemsep=0.1em,topsep=0.3em,leftmargin=1.4em}
\renewcommand{\arraystretch}{1.15}
\setlength{\tabcolsep}{5pt}
\setlength{\LTpre}{0.8em}
\setlength{\LTpost}{0.4em}

% Legendas no padrão Elsevier: \enquote{Fig. 1.} seguido do texto, alinhado à esquerda;
% \enquote{Tabela 1} numa linha e o título na linha seguinte
\DeclareCaptionLabelSeparator{travessao}{\ — }
% `babel` guarda os nomes por idioma, e o nome da opção varia entre os
% documentos. Entra em cada um que estiver carregado.
\@for\pb@idioma:=brazilian,brazil,portuguese,portuges,american,english\do{%
  \@ifundefined{captions\pb@idioma}{}{%
    \expandafter\addto\csname captions\pb@idioma\endcsname{\renewcommand{\figurename}{Fig.}}}}
\captionsetup{font=footnotesize,labelfont=bf,justification=justified,
              singlelinecheck=false,skip=5pt}
\captionsetup[figure]{position=bottom,labelsep=period}
\captionsetup[table]{position=top,labelsep=newline}
% DIFERENÇA: a origem fixa todo flutuante onde aparece (`H`). Aqui ele
% tenta o lugar e, se não couber no resto da página, sobe para a seguinte
% e o texto continua: com `H`, a tabela que não cabia deixava o fim da
% página em branco.
\floatplacement{table}{!htbp}
\floatplacement{figure}{!htbp}
\renewcommand{\topfraction}{0.9}
\renewcommand{\bottomfraction}{0.8}
\renewcommand{\textfraction}{0.08}
\renewcommand{\floatpagefraction}{0.75}

\newcolumntype{Y}{>{\raggedright\arraybackslash}X}
\newcolumntype{L}[1]{>{\raggedright\arraybackslash}p{#1}}
\newcolumntype{C}[1]{>{\centering\arraybackslash}p{#1}}
\newcolumntype{R}[1]{>{\raggedleft\arraybackslash}p{#1}}

% ---------------------------------------------------------------------
%  Títulos de seção (ISO 2145)
% ---------------------------------------------------------------------
% DIFERENÇA: título de manual. Seção e subseção avançam 14 mm para a
% esquerda do texto e ocupam a largura útil; a seção leva um fio acima.
% Número sem ponto final, três níveis numerados. Os três níveis têm o
% corpo do texto ou pouco mais: o nível se lê pelo fio, pelo recuo e
% pelo itálico do terceiro, e não pelo tamanho.
\setcounter{secnumdepth}{3}
\titleformat{\section}[block]
  {\titlerule[0.8pt]\vspace{0.5ex}\normalfont\large\bfseries\raggedright}{\thesection}{0.8em}{}
\titleformat{\subsection}{\normalfont\normalsize\bfseries\raggedright}{\thesubsection}{0.8em}{}
\titleformat{\subsubsection}{\normalfont\normalsize\bfseries\itshape\raggedright}{\thesubsubsection}{0.8em}{}
\titlespacing*{\section}{-\margemmanual}{1.5em plus 0.4em minus 0.2em}{0.5em plus 0.1em}
\titlespacing*{\subsection}{-\margemmanual}{1.1em plus 0.3em minus 0.2em}{0.3em plus 0.1em}
\titlespacing*{\subsubsection}{0pt}{0.9em plus 0.2em minus 0.2em}{0.2em}

% ---------------------------------------------------------------------
%  Identificação, cabeçalho e rodapé
% ---------------------------------------------------------------------
% Todo campo é do documento, por `\renewcommand` depois deste arquivo.
% Sem isso: o código é o nome do arquivo, a data é a da compilação em
% ISO 8601, e versão e responsável ficam vazios e somem.
\providecommand{\doctitulo}{}
\providecommand{\docsubtitulo}{}
\providecommand{\doccodigo}{\jobname}
\providecommand{\docversao}{}
\providecommand{\docresponsavel}{}
\providecommand{\docdata}{\the\year-\two@digits\month-\two@digits\day}

% Cabeçalho: título à esquerda, seção corrente à direita. Rodapé: código,
% versão e data à esquerda, "página X de Y" à direita. Os dois ocupam a
% largura útil, margem de título incluída.
\newcommand{\pb@identificacao}{%
  \texttt{\doccodigo}%
  \ifx\docversao\@empty\else\enspace\textperiodcentered\enspace versão \docversao\fi
  \enspace\textperiodcentered\enspace\docdata}
\pagestyle{fancy}
\fancyhf{}
\fancyhfoffset[L]{\margemmanual}
\fancyhead[L]{\footnotesize\bfseries\doctitulo}
\fancyhead[R]{\footnotesize\itshape\nouppercase{\rightmark}}
\fancyfoot[L]{\footnotesize\pb@identificacao}
\fancyfoot[R]{\footnotesize\thepage\ de \pageref*{LastPage}}
\renewcommand{\headrulewidth}{0.4pt}
\renewcommand{\footrulewidth}{0.4pt}
\fancypagestyle{plain}{}
% Página de rosto: sem cabeçalho, porque o título já está nela.
\fancypagestyle{rosto}{\fancyhead{}\renewcommand{\headrulewidth}{0pt}}

% ---------------------------------------------------------------------
%  Marca
% ---------------------------------------------------------------------
% DIFERENÇA: a origem não tem marca. Aqui todo documento abre com o
% logotipo Solara, versão para fundo claro. O caminho depende de onde se
% compila: `report/` ou `docs/product/`.
\IfFileExists{../docs/img/solara-logo-on-light.png}%
  {\newcommand{\marcaarquivo}{../docs/img/solara-logo-on-light.png}}%
  {\newcommand{\marcaarquivo}{../img/solara-logo-on-light.png}}
% \marcadoc[largura]: o logotipo tem proporção 7:1, então 30 mm dão
% 4,3 mm de altura.
\newcommand{\marcadoc}[1][30mm]{\includegraphics[width=#1]{\marcaarquivo}}

% ---------------------------------------------------------------------
%  Rosto e histórico de revisões
% ---------------------------------------------------------------------
% \rosto: marca, título, subtítulo e a tabela de identificação. Lê os
% campos `\doc...`; linha de campo vazio não sai.
\newcommand{\pb@linha}[2]{\ifx#2\@empty\else #1 & #2 \\\fi}
\newcommand{\rosto}{%
  \thispagestyle{rosto}%
  \begin{adjustwidth}{-\margemmanual}{0pt}
    \raggedright
    \marcadoc[45mm]\par\vspace{16mm}
    {\fontsize{24}{29}\selectfont\bfseries\doctitulo\par}
    \ifx\docsubtitulo\@empty\else
      \vspace{3mm}{\Large\color{cinza}\docsubtitulo\par}\fi
    \vspace{9mm}{\hrule height 0.8pt}\vspace{4mm}
    {\small\renewcommand{\arraystretch}{1.25}%
     \begin{tabular}{@{}>{\bfseries}p{28mm}p{110mm}@{}}
       Documento & \texttt{\doccodigo} \\
       \pb@linha{Versão}{\docversao}%
       Data & \docdata \\
       \pb@linha{Responsável}{\docresponsavel}%
     \end{tabular}\par}
  \end{adjustwidth}
  \vspace{10mm}}

% Histórico de revisões: uma linha por versão, a mais recente em cima.
%   \begin{historico}
%     0.2 & 2026-10-03 & O que mudou \\
%   \end{historico}
\newenvironment{historico}
  {\par\medskip\noindent\small
   \begin{tabular}{@{}p{16mm}p{24mm}p{\dimexpr\linewidth-40mm-4\tabcolsep\relax}@{}}
   \toprule \textbf{Versão} & \textbf{Data} & \textbf{Alteração} \\ \midrule}
  {\bottomrule\end{tabular}\par\medskip}

% ---------------------------------------------------------------------
%  Avisos (ANSI Z535.6) e nota
% ---------------------------------------------------------------------
% Quatro palavras de sinal, em hierarquia fixa. As três primeiras são de
% dano a pessoa; AVISO é de dano a equipamento ou a dado. Cada mensagem
% diz o risco, como evitar e a consequência. NOTA não é aviso: é
% informação de apoio.
\definecolor{sinalperigo}{HTML}{C8102E}
\definecolor{sinalatencao}{HTML}{E87722}
\definecolor{sinalcuidado}{HTML}{F2C200}
\definecolor{sinalaviso}{HTML}{00539B}
\definecolor{sinalnota}{HTML}{5F5F5F}
% #1 palavra, #2 cor da faixa, #3 cor da letra sobre a faixa
\newtcolorbox{pb@mensagem}[3]{enhanced,breakable,sharp corners,
  boxrule=0pt,frame hidden,borderline west={3pt}{0pt}{#2},
  colback=#2!7,left=9pt,right=7pt,top=5pt,bottom=5pt,
  before skip=0.9em,after skip=0.9em,fontupper=\small,
  before upper={\colorbox{#2}{\textcolor{#3}{\footnotesize\bfseries #1}}\enspace}}
\newenvironment{perigo}{\csname pb@mensagem\endcsname{PERIGO}{sinalperigo}{white}}{\csname endpb@mensagem\endcsname}
\newenvironment{atencao}{\csname pb@mensagem\endcsname{ATENÇÃO}{sinalatencao}{black}}{\csname endpb@mensagem\endcsname}
\newenvironment{cuidado}{\csname pb@mensagem\endcsname{CUIDADO}{sinalcuidado}{black}}{\csname endpb@mensagem\endcsname}
\newenvironment{aviso}{\csname pb@mensagem\endcsname{AVISO}{sinalaviso}{white}}{\csname endpb@mensagem\endcsname}
\newenvironment{nota}{\csname pb@mensagem\endcsname{NOTA}{sinalnota}{white}}{\csname endpb@mensagem\endcsname}

% ---------------------------------------------------------------------
%  Procedimento
% ---------------------------------------------------------------------
% Passos numerados, uma ação por passo, no imperativo. A pré-condição vem
% antes da lista e o resultado esperado depois.
\newlist{passos}{enumerate}{2}
\setlist[passos,1]{label=\textbf{\arabic*.},leftmargin=2.2em,itemsep=0.45em,topsep=0.6em}
\setlist[passos,2]{label=\alph*),leftmargin=1.8em,itemsep=0.2em,topsep=0.2em}
\newcommand{\precondicao}[1]{\par\smallskip\noindent\textbf{Antes de começar.}\enspace#1\par}
\newcommand{\resultado}[1]{\par\smallskip\noindent\textbf{Resultado.}\enspace#1\par}

% Bloco na largura útil inteira, margem de título incluída: para tabela
% ou figura que não cabe nos 152 mm do texto.
\newenvironment{larga}{\begin{adjustwidth}{-\margemmanual}{0pt}}{\end{adjustwidth}}
% O mesmo dentro de `table` ou `figure`, onde `adjustwidth` não alcança:
% uma caixa na largura útil, puxada para a margem. Não quebra página.
% Dentro dela a largura é `\linewidth`, e não `\textwidth`.
\newenvironment{tabelalarga}
  {\noindent\hspace*{-\margemmanual}%
   \begin{minipage}{\dimexpr\linewidth+\margemmanual\relax}\centering}
  {\end{minipage}}

% ---------------------------------------------------------------------
%  Campos e figuras
% ---------------------------------------------------------------------
\definecolor{cinza}{gray}{0.40}

% \figuradoc[largura]{caminho}; sem o arquivo, uma caixa marca a pendência
\newcommand{\figuradoc}[2][\linewidth]{%
  \IfFileExists{#2}%
    {\includegraphics[width=#1,height=0.80\textheight,keepaspectratio]{#2}}%
    {\fbox{\parbox[c][45mm][c]{\dimexpr#1-2\fboxsep-2\fboxrule\relax}%
      {\centering\small\color{cinza}\itshape[figura em preparação]}}}}

\pbrestauraarroba

%
%  FORMA DE MANUAL TÉCNICO. Fonte, números, legendas e pacotes vêm da
%  origem. A forma da página não: é de manual de referência, e não de
%  artigo. Coluna única, títulos recuados para a esquerda do texto,
%  rodapé com identificação e "página X de Y" em toda página, avisos com
%  palavra de sinal (ANSI Z535.6), passos numerados e histórico de
%  revisões (IEC/IEEE 82079-1). O desenho de página segue a classe
%  `refman` do CTAN.
%
%  O documento carrega `babel` ANTES deste arquivo. Documento deitado, que
%  existe para caber tabela larga, carrega também `geometry` antes; aí a
%  margem de manual não entra e os títulos não avançam.
%
%  Motor: pdflatex. Charis SIL, newtxmath e Bera Mono entram como Type 1
%  em T1, e a matemática do newtxmath não passa pelo `unicode-math`.
% =====================================================================
\RequirePackage{iftex}
\RequirePDFTeX
% Este arquivo entra tanto de um `.tex` quanto de dentro do
% `estilo-reporte.sty`, onde o `@` ja e letra. O codigo de categoria e
% devolvido como estava, e nao forcado a "outro" no fim.
\edef\pbrestauraarroba{\catcode`\noexpand\@=\the\catcode`\@\relax}
\makeatletter

% ---------------------------------------------------------------------
%  Fontes, idioma e microtipografia
% ---------------------------------------------------------------------
\usepackage[T1]{fontenc}
\usepackage[utf8]{inputenc}
% DIFERENÇA: o idioma é do documento, porque há documento em inglês.
\@ifpackageloaded{babel}{}{\usepackage[brazilian]{babel}}
\usepackage{CharisSIL}
\usepackage[charter]{newtxmath}
\usepackage[scaled=0.83]{beramono}
\usepackage[protrusion=true,expansion=true,tracking=true]{microtype}
\SetTracking{encoding=*,shape=sc}{40}
% DIFERENÇA: a origem não usa sem serifa. Os documentos daqui pedem
% `\sffamily` em título, tabela e nota; sem esta linha ela cairia na
% Computer Modern Sans, que não pertence ao conjunto. Uma família só.
\renewcommand{\sfdefault}{\rmdefault}

% DIFERENÇA: mancha de manual, e não a coluna de 468 pt da origem. A4 com
% 22 mm de margem de cada lado, 166 mm úteis. O texto corre em 152 mm e
% os 14 mm da esquerda são o recuo dos títulos: o bastante para o título
% sair do alinhamento do texto, sem deixar coluna vazia ao lado dele. A
% primeira versão reservava 40 mm, e a margem ficava em branco em toda
% página que não fosse título; o documento crescia um terço só por isso.
% `\margemmanual` é quanto os títulos, o cabeçalho e o rodapé avançam
% para a esquerda; vale zero quando o documento traz mancha própria.
% A altura é de 253 mm, a mesma da mancha anterior, porque as tabelas
% deitadas foram dimensionadas para ela.
\newlength{\margemmanual}
\@ifpackageloaded{geometry}{\setlength{\margemmanual}{0pt}}{%
  \setlength{\margemmanual}{14mm}%
  \usepackage[a4paper,left=36mm,right=22mm,top=20mm,bottom=24mm,
              headheight=14pt,headsep=7mm,footskip=12mm]{geometry}}

\usepackage{amsmath}
\usepackage[dvipsnames,table]{xcolor}
\usepackage{graphicx}
\usepackage{tikz}
\usepackage{booktabs}
\usepackage{tabularx}
\usepackage{array}
\usepackage{xltabular}
\usepackage{float}
\usepackage{pdflscape}
\usepackage{rotating}
% Página com o conteúdo girado, mas exibida em retrato no leitor de PDF
\newenvironment{paisagemretrato}{\begingroup\def\PLS@AddRotate##1{}\landscape}{\endlandscape\endgroup}
\usepackage{enumitem}
\usepackage{changepage}
\usepackage[most]{tcolorbox}
\usepackage{caption}
\usepackage{titlesec}
\usepackage{fancyhdr}
\usepackage{lastpage}
\usepackage{csquotes}
\usepackage{siunitx}
% DIFERENÇA: sem `biblatex`. Nenhum documento daqui tem `.bib`; as
% referências são listas escritas à mão.
\usepackage{xurl}
\usepackage[hidelinks,pdfusetitle]{hyperref}

% ---------------------------------------------------------------------
%  Números e unidades (ISO 80000-1)
% ---------------------------------------------------------------------
\sisetup{
  output-decimal-marker = {,},
  group-separator       = {\,},
  group-minimum-digits  = 4,
  exponent-product      = \cdot,
  uncertainty-mode      = separate,
  range-phrase          = {\,a\,},
  list-pair-separator   = {\ e\ },
  list-final-separator  = {\ e\ },
  per-mode              = symbol,
}

% ---------------------------------------------------------------------
%  Parágrafo, listas e flutuantes
% ---------------------------------------------------------------------
% DIFERENÇA: parágrafo em bloco, separado por espaço e sem recuo, que é
% o de manual; a origem recua 1 em.
\setlength{\parindent}{0pt}
\setlength{\parskip}{0.45em plus 0.1em}
\linespread{1.04}
\setlist{itemsep=0.1em,topsep=0.3em,leftmargin=1.4em}
\renewcommand{\arraystretch}{1.15}
\setlength{\tabcolsep}{5pt}
\setlength{\LTpre}{0.8em}
\setlength{\LTpost}{0.4em}

% Legendas no padrão Elsevier: \enquote{Fig. 1.} seguido do texto, alinhado à esquerda;
% \enquote{Tabela 1} numa linha e o título na linha seguinte
\DeclareCaptionLabelSeparator{travessao}{\ — }
% `babel` guarda os nomes por idioma, e o nome da opção varia entre os
% documentos. Entra em cada um que estiver carregado.
\@for\pb@idioma:=brazilian,brazil,portuguese,portuges,american,english\do{%
  \@ifundefined{captions\pb@idioma}{}{%
    \expandafter\addto\csname captions\pb@idioma\endcsname{\renewcommand{\figurename}{Fig.}}}}
\captionsetup{font=footnotesize,labelfont=bf,justification=justified,
              singlelinecheck=false,skip=5pt}
\captionsetup[figure]{position=bottom,labelsep=period}
\captionsetup[table]{position=top,labelsep=newline}
% DIFERENÇA: a origem fixa todo flutuante onde aparece (`H`). Aqui ele
% tenta o lugar e, se não couber no resto da página, sobe para a seguinte
% e o texto continua: com `H`, a tabela que não cabia deixava o fim da
% página em branco.
\floatplacement{table}{!htbp}
\floatplacement{figure}{!htbp}
\renewcommand{\topfraction}{0.9}
\renewcommand{\bottomfraction}{0.8}
\renewcommand{\textfraction}{0.08}
\renewcommand{\floatpagefraction}{0.75}

\newcolumntype{Y}{>{\raggedright\arraybackslash}X}
\newcolumntype{L}[1]{>{\raggedright\arraybackslash}p{#1}}
\newcolumntype{C}[1]{>{\centering\arraybackslash}p{#1}}
\newcolumntype{R}[1]{>{\raggedleft\arraybackslash}p{#1}}

% ---------------------------------------------------------------------
%  Títulos de seção (ISO 2145)
% ---------------------------------------------------------------------
% DIFERENÇA: título de manual. Seção e subseção avançam 14 mm para a
% esquerda do texto e ocupam a largura útil; a seção leva um fio acima.
% Número sem ponto final, três níveis numerados. Os três níveis têm o
% corpo do texto ou pouco mais: o nível se lê pelo fio, pelo recuo e
% pelo itálico do terceiro, e não pelo tamanho.
\setcounter{secnumdepth}{3}
\titleformat{\section}[block]
  {\titlerule[0.8pt]\vspace{0.5ex}\normalfont\large\bfseries\raggedright}{\thesection}{0.8em}{}
\titleformat{\subsection}{\normalfont\normalsize\bfseries\raggedright}{\thesubsection}{0.8em}{}
\titleformat{\subsubsection}{\normalfont\normalsize\bfseries\itshape\raggedright}{\thesubsubsection}{0.8em}{}
\titlespacing*{\section}{-\margemmanual}{1.5em plus 0.4em minus 0.2em}{0.5em plus 0.1em}
\titlespacing*{\subsection}{-\margemmanual}{1.1em plus 0.3em minus 0.2em}{0.3em plus 0.1em}
\titlespacing*{\subsubsection}{0pt}{0.9em plus 0.2em minus 0.2em}{0.2em}

% ---------------------------------------------------------------------
%  Identificação, cabeçalho e rodapé
% ---------------------------------------------------------------------
% Todo campo é do documento, por `\renewcommand` depois deste arquivo.
% Sem isso: o código é o nome do arquivo, a data é a da compilação em
% ISO 8601, e versão e responsável ficam vazios e somem.
\providecommand{\doctitulo}{}
\providecommand{\docsubtitulo}{}
\providecommand{\doccodigo}{\jobname}
\providecommand{\docversao}{}
\providecommand{\docresponsavel}{}
\providecommand{\docdata}{\the\year-\two@digits\month-\two@digits\day}

% Cabeçalho: título à esquerda, seção corrente à direita. Rodapé: código,
% versão e data à esquerda, "página X de Y" à direita. Os dois ocupam a
% largura útil, margem de título incluída.
\newcommand{\pb@identificacao}{%
  \texttt{\doccodigo}%
  \ifx\docversao\@empty\else\enspace\textperiodcentered\enspace versão \docversao\fi
  \enspace\textperiodcentered\enspace\docdata}
\pagestyle{fancy}
\fancyhf{}
\fancyhfoffset[L]{\margemmanual}
\fancyhead[L]{\footnotesize\bfseries\doctitulo}
\fancyhead[R]{\footnotesize\itshape\nouppercase{\rightmark}}
\fancyfoot[L]{\footnotesize\pb@identificacao}
\fancyfoot[R]{\footnotesize\thepage\ de \pageref*{LastPage}}
\renewcommand{\headrulewidth}{0.4pt}
\renewcommand{\footrulewidth}{0.4pt}
\fancypagestyle{plain}{}
% Página de rosto: sem cabeçalho, porque o título já está nela.
\fancypagestyle{rosto}{\fancyhead{}\renewcommand{\headrulewidth}{0pt}}

% ---------------------------------------------------------------------
%  Marca
% ---------------------------------------------------------------------
% DIFERENÇA: a origem não tem marca. Aqui todo documento abre com o
% logotipo Solara, versão para fundo claro. O caminho depende de onde se
% compila: `report/` ou `docs/product/`.
\IfFileExists{../docs/img/solara-logo-on-light.png}%
  {\newcommand{\marcaarquivo}{../docs/img/solara-logo-on-light.png}}%
  {\newcommand{\marcaarquivo}{../img/solara-logo-on-light.png}}
% \marcadoc[largura]: o logotipo tem proporção 7:1, então 30 mm dão
% 4,3 mm de altura.
\newcommand{\marcadoc}[1][30mm]{\includegraphics[width=#1]{\marcaarquivo}}

% ---------------------------------------------------------------------
%  Rosto e histórico de revisões
% ---------------------------------------------------------------------
% \rosto: marca, título, subtítulo e a tabela de identificação. Lê os
% campos `\doc...`; linha de campo vazio não sai.
\newcommand{\pb@linha}[2]{\ifx#2\@empty\else #1 & #2 \\\fi}
\newcommand{\rosto}{%
  \thispagestyle{rosto}%
  \begin{adjustwidth}{-\margemmanual}{0pt}
    \raggedright
    \marcadoc[45mm]\par\vspace{16mm}
    {\fontsize{24}{29}\selectfont\bfseries\doctitulo\par}
    \ifx\docsubtitulo\@empty\else
      \vspace{3mm}{\Large\color{cinza}\docsubtitulo\par}\fi
    \vspace{9mm}{\hrule height 0.8pt}\vspace{4mm}
    {\small\renewcommand{\arraystretch}{1.25}%
     \begin{tabular}{@{}>{\bfseries}p{28mm}p{110mm}@{}}
       Documento & \texttt{\doccodigo} \\
       \pb@linha{Versão}{\docversao}%
       Data & \docdata \\
       \pb@linha{Responsável}{\docresponsavel}%
     \end{tabular}\par}
  \end{adjustwidth}
  \vspace{10mm}}

% Histórico de revisões: uma linha por versão, a mais recente em cima.
%   \begin{historico}
%     0.2 & 2026-10-03 & O que mudou \\
%   \end{historico}
\newenvironment{historico}
  {\par\medskip\noindent\small
   \begin{tabular}{@{}p{16mm}p{24mm}p{\dimexpr\linewidth-40mm-4\tabcolsep\relax}@{}}
   \toprule \textbf{Versão} & \textbf{Data} & \textbf{Alteração} \\ \midrule}
  {\bottomrule\end{tabular}\par\medskip}

% ---------------------------------------------------------------------
%  Avisos (ANSI Z535.6) e nota
% ---------------------------------------------------------------------
% Quatro palavras de sinal, em hierarquia fixa. As três primeiras são de
% dano a pessoa; AVISO é de dano a equipamento ou a dado. Cada mensagem
% diz o risco, como evitar e a consequência. NOTA não é aviso: é
% informação de apoio.
\definecolor{sinalperigo}{HTML}{C8102E}
\definecolor{sinalatencao}{HTML}{E87722}
\definecolor{sinalcuidado}{HTML}{F2C200}
\definecolor{sinalaviso}{HTML}{00539B}
\definecolor{sinalnota}{HTML}{5F5F5F}
% #1 palavra, #2 cor da faixa, #3 cor da letra sobre a faixa
\newtcolorbox{pb@mensagem}[3]{enhanced,breakable,sharp corners,
  boxrule=0pt,frame hidden,borderline west={3pt}{0pt}{#2},
  colback=#2!7,left=9pt,right=7pt,top=5pt,bottom=5pt,
  before skip=0.9em,after skip=0.9em,fontupper=\small,
  before upper={\colorbox{#2}{\textcolor{#3}{\footnotesize\bfseries #1}}\enspace}}
\newenvironment{perigo}{\csname pb@mensagem\endcsname{PERIGO}{sinalperigo}{white}}{\csname endpb@mensagem\endcsname}
\newenvironment{atencao}{\csname pb@mensagem\endcsname{ATENÇÃO}{sinalatencao}{black}}{\csname endpb@mensagem\endcsname}
\newenvironment{cuidado}{\csname pb@mensagem\endcsname{CUIDADO}{sinalcuidado}{black}}{\csname endpb@mensagem\endcsname}
\newenvironment{aviso}{\csname pb@mensagem\endcsname{AVISO}{sinalaviso}{white}}{\csname endpb@mensagem\endcsname}
\newenvironment{nota}{\csname pb@mensagem\endcsname{NOTA}{sinalnota}{white}}{\csname endpb@mensagem\endcsname}

% ---------------------------------------------------------------------
%  Procedimento
% ---------------------------------------------------------------------
% Passos numerados, uma ação por passo, no imperativo. A pré-condição vem
% antes da lista e o resultado esperado depois.
\newlist{passos}{enumerate}{2}
\setlist[passos,1]{label=\textbf{\arabic*.},leftmargin=2.2em,itemsep=0.45em,topsep=0.6em}
\setlist[passos,2]{label=\alph*),leftmargin=1.8em,itemsep=0.2em,topsep=0.2em}
\newcommand{\precondicao}[1]{\par\smallskip\noindent\textbf{Antes de começar.}\enspace#1\par}
\newcommand{\resultado}[1]{\par\smallskip\noindent\textbf{Resultado.}\enspace#1\par}

% Bloco na largura útil inteira, margem de título incluída: para tabela
% ou figura que não cabe nos 152 mm do texto.
\newenvironment{larga}{\begin{adjustwidth}{-\margemmanual}{0pt}}{\end{adjustwidth}}
% O mesmo dentro de `table` ou `figure`, onde `adjustwidth` não alcança:
% uma caixa na largura útil, puxada para a margem. Não quebra página.
% Dentro dela a largura é `\linewidth`, e não `\textwidth`.
\newenvironment{tabelalarga}
  {\noindent\hspace*{-\margemmanual}%
   \begin{minipage}{\dimexpr\linewidth+\margemmanual\relax}\centering}
  {\end{minipage}}

% ---------------------------------------------------------------------
%  Campos e figuras
% ---------------------------------------------------------------------
\definecolor{cinza}{gray}{0.40}

% \figuradoc[largura]{caminho}; sem o arquivo, uma caixa marca a pendência
\newcommand{\figuradoc}[2][\linewidth]{%
  \IfFileExists{#2}%
    {\includegraphics[width=#1,height=0.80\textheight,keepaspectratio]{#2}}%
    {\fbox{\parbox[c][45mm][c]{\dimexpr#1-2\fboxsep-2\fboxrule\relax}%
      {\centering\small\color{cinza}\itshape[figura em preparação]}}}}

\pbrestauraarroba

%
%  FORMA DE MANUAL TÉCNICO. Fonte, números, legendas e pacotes vêm da
%  origem. A forma da página não: é de manual de referência, e não de
%  artigo. Coluna única, títulos recuados para a esquerda do texto,
%  rodapé com identificação e "página X de Y" em toda página, avisos com
%  palavra de sinal (ANSI Z535.6), passos numerados e histórico de
%  revisões (IEC/IEEE 82079-1). O desenho de página segue a classe
%  `refman` do CTAN.
%
%  O documento carrega `babel` ANTES deste arquivo. Documento deitado, que
%  existe para caber tabela larga, carrega também `geometry` antes; aí a
%  margem de manual não entra e os títulos não avançam.
%
%  Motor: pdflatex. Charis SIL, newtxmath e Bera Mono entram como Type 1
%  em T1, e a matemática do newtxmath não passa pelo `unicode-math`.
% =====================================================================
\RequirePackage{iftex}
\RequirePDFTeX
% Este arquivo entra tanto de um `.tex` quanto de dentro do
% `estilo-reporte.sty`, onde o `@` ja e letra. O codigo de categoria e
% devolvido como estava, e nao forcado a "outro" no fim.
\edef\pbrestauraarroba{\catcode`\noexpand\@=\the\catcode`\@\relax}
\makeatletter

% ---------------------------------------------------------------------
%  Fontes, idioma e microtipografia
% ---------------------------------------------------------------------
\usepackage[T1]{fontenc}
\usepackage[utf8]{inputenc}
% DIFERENÇA: o idioma é do documento, porque há documento em inglês.
\@ifpackageloaded{babel}{}{\usepackage[brazilian]{babel}}
\usepackage{CharisSIL}
\usepackage[charter]{newtxmath}
\usepackage[scaled=0.83]{beramono}
\usepackage[protrusion=true,expansion=true,tracking=true]{microtype}
\SetTracking{encoding=*,shape=sc}{40}
% DIFERENÇA: a origem não usa sem serifa. Os documentos daqui pedem
% `\sffamily` em título, tabela e nota; sem esta linha ela cairia na
% Computer Modern Sans, que não pertence ao conjunto. Uma família só.
\renewcommand{\sfdefault}{\rmdefault}

% DIFERENÇA: mancha de manual, e não a coluna de 468 pt da origem. A4 com
% 22 mm de margem de cada lado, 166 mm úteis. O texto corre em 152 mm e
% os 14 mm da esquerda são o recuo dos títulos: o bastante para o título
% sair do alinhamento do texto, sem deixar coluna vazia ao lado dele. A
% primeira versão reservava 40 mm, e a margem ficava em branco em toda
% página que não fosse título; o documento crescia um terço só por isso.
% `\margemmanual` é quanto os títulos, o cabeçalho e o rodapé avançam
% para a esquerda; vale zero quando o documento traz mancha própria.
% A altura é de 253 mm, a mesma da mancha anterior, porque as tabelas
% deitadas foram dimensionadas para ela.
\newlength{\margemmanual}
\@ifpackageloaded{geometry}{\setlength{\margemmanual}{0pt}}{%
  \setlength{\margemmanual}{14mm}%
  \usepackage[a4paper,left=36mm,right=22mm,top=20mm,bottom=24mm,
              headheight=14pt,headsep=7mm,footskip=12mm]{geometry}}

\usepackage{amsmath}
\usepackage[dvipsnames,table]{xcolor}
\usepackage{graphicx}
\usepackage{tikz}
\usepackage{booktabs}
\usepackage{tabularx}
\usepackage{array}
\usepackage{xltabular}
\usepackage{float}
\usepackage{pdflscape}
\usepackage{rotating}
% Página com o conteúdo girado, mas exibida em retrato no leitor de PDF
\newenvironment{paisagemretrato}{\begingroup\def\PLS@AddRotate##1{}\landscape}{\endlandscape\endgroup}
\usepackage{enumitem}
\usepackage{changepage}
\usepackage[most]{tcolorbox}
\usepackage{caption}
\usepackage{titlesec}
\usepackage{fancyhdr}
\usepackage{lastpage}
\usepackage{csquotes}
\usepackage{siunitx}
% DIFERENÇA: sem `biblatex`. Nenhum documento daqui tem `.bib`; as
% referências são listas escritas à mão.
\usepackage{xurl}
\usepackage[hidelinks,pdfusetitle]{hyperref}

% ---------------------------------------------------------------------
%  Números e unidades (ISO 80000-1)
% ---------------------------------------------------------------------
\sisetup{
  output-decimal-marker = {,},
  group-separator       = {\,},
  group-minimum-digits  = 4,
  exponent-product      = \cdot,
  uncertainty-mode      = separate,
  range-phrase          = {\,a\,},
  list-pair-separator   = {\ e\ },
  list-final-separator  = {\ e\ },
  per-mode              = symbol,
}

% ---------------------------------------------------------------------
%  Parágrafo, listas e flutuantes
% ---------------------------------------------------------------------
% DIFERENÇA: parágrafo em bloco, separado por espaço e sem recuo, que é
% o de manual; a origem recua 1 em.
\setlength{\parindent}{0pt}
\setlength{\parskip}{0.45em plus 0.1em}
\linespread{1.04}
\setlist{itemsep=0.1em,topsep=0.3em,leftmargin=1.4em}
\renewcommand{\arraystretch}{1.15}
\setlength{\tabcolsep}{5pt}
\setlength{\LTpre}{0.8em}
\setlength{\LTpost}{0.4em}

% Legendas no padrão Elsevier: \enquote{Fig. 1.} seguido do texto, alinhado à esquerda;
% \enquote{Tabela 1} numa linha e o título na linha seguinte
\DeclareCaptionLabelSeparator{travessao}{\ — }
% `babel` guarda os nomes por idioma, e o nome da opção varia entre os
% documentos. Entra em cada um que estiver carregado.
\@for\pb@idioma:=brazilian,brazil,portuguese,portuges,american,english\do{%
  \@ifundefined{captions\pb@idioma}{}{%
    \expandafter\addto\csname captions\pb@idioma\endcsname{\renewcommand{\figurename}{Fig.}}}}
\captionsetup{font=footnotesize,labelfont=bf,justification=justified,
              singlelinecheck=false,skip=5pt}
\captionsetup[figure]{position=bottom,labelsep=period}
\captionsetup[table]{position=top,labelsep=newline}
% DIFERENÇA: a origem fixa todo flutuante onde aparece (`H`). Aqui ele
% tenta o lugar e, se não couber no resto da página, sobe para a seguinte
% e o texto continua: com `H`, a tabela que não cabia deixava o fim da
% página em branco.
\floatplacement{table}{!htbp}
\floatplacement{figure}{!htbp}
\renewcommand{\topfraction}{0.9}
\renewcommand{\bottomfraction}{0.8}
\renewcommand{\textfraction}{0.08}
\renewcommand{\floatpagefraction}{0.75}

\newcolumntype{Y}{>{\raggedright\arraybackslash}X}
\newcolumntype{L}[1]{>{\raggedright\arraybackslash}p{#1}}
\newcolumntype{C}[1]{>{\centering\arraybackslash}p{#1}}
\newcolumntype{R}[1]{>{\raggedleft\arraybackslash}p{#1}}

% ---------------------------------------------------------------------
%  Títulos de seção (ISO 2145)
% ---------------------------------------------------------------------
% DIFERENÇA: título de manual. Seção e subseção avançam 14 mm para a
% esquerda do texto e ocupam a largura útil; a seção leva um fio acima.
% Número sem ponto final, três níveis numerados. Os três níveis têm o
% corpo do texto ou pouco mais: o nível se lê pelo fio, pelo recuo e
% pelo itálico do terceiro, e não pelo tamanho.
\setcounter{secnumdepth}{3}
\titleformat{\section}[block]
  {\titlerule[0.8pt]\vspace{0.5ex}\normalfont\large\bfseries\raggedright}{\thesection}{0.8em}{}
\titleformat{\subsection}{\normalfont\normalsize\bfseries\raggedright}{\thesubsection}{0.8em}{}
\titleformat{\subsubsection}{\normalfont\normalsize\bfseries\itshape\raggedright}{\thesubsubsection}{0.8em}{}
\titlespacing*{\section}{-\margemmanual}{1.5em plus 0.4em minus 0.2em}{0.5em plus 0.1em}
\titlespacing*{\subsection}{-\margemmanual}{1.1em plus 0.3em minus 0.2em}{0.3em plus 0.1em}
\titlespacing*{\subsubsection}{0pt}{0.9em plus 0.2em minus 0.2em}{0.2em}

% ---------------------------------------------------------------------
%  Identificação, cabeçalho e rodapé
% ---------------------------------------------------------------------
% Todo campo é do documento, por `\renewcommand` depois deste arquivo.
% Sem isso: o código é o nome do arquivo, a data é a da compilação em
% ISO 8601, e versão e responsável ficam vazios e somem.
\providecommand{\doctitulo}{}
\providecommand{\docsubtitulo}{}
\providecommand{\doccodigo}{\jobname}
\providecommand{\docversao}{}
\providecommand{\docresponsavel}{}
\providecommand{\docdata}{\the\year-\two@digits\month-\two@digits\day}

% Cabeçalho: título à esquerda, seção corrente à direita. Rodapé: código,
% versão e data à esquerda, "página X de Y" à direita. Os dois ocupam a
% largura útil, margem de título incluída.
\newcommand{\pb@identificacao}{%
  \texttt{\doccodigo}%
  \ifx\docversao\@empty\else\enspace\textperiodcentered\enspace versão \docversao\fi
  \enspace\textperiodcentered\enspace\docdata}
\pagestyle{fancy}
\fancyhf{}
\fancyhfoffset[L]{\margemmanual}
\fancyhead[L]{\footnotesize\bfseries\doctitulo}
\fancyhead[R]{\footnotesize\itshape\nouppercase{\rightmark}}
\fancyfoot[L]{\footnotesize\pb@identificacao}
\fancyfoot[R]{\footnotesize\thepage\ de \pageref*{LastPage}}
\renewcommand{\headrulewidth}{0.4pt}
\renewcommand{\footrulewidth}{0.4pt}
\fancypagestyle{plain}{}
% Página de rosto: sem cabeçalho, porque o título já está nela.
\fancypagestyle{rosto}{\fancyhead{}\renewcommand{\headrulewidth}{0pt}}

% ---------------------------------------------------------------------
%  Marca
% ---------------------------------------------------------------------
% DIFERENÇA: a origem não tem marca. Aqui todo documento abre com o
% logotipo Solara, versão para fundo claro. O caminho depende de onde se
% compila: `report/` ou `docs/product/`.
\IfFileExists{../docs/img/solara-logo-on-light.png}%
  {\newcommand{\marcaarquivo}{../docs/img/solara-logo-on-light.png}}%
  {\newcommand{\marcaarquivo}{../img/solara-logo-on-light.png}}
% \marcadoc[largura]: o logotipo tem proporção 7:1, então 30 mm dão
% 4,3 mm de altura.
\newcommand{\marcadoc}[1][30mm]{\includegraphics[width=#1]{\marcaarquivo}}

% ---------------------------------------------------------------------
%  Rosto e histórico de revisões
% ---------------------------------------------------------------------
% \rosto: marca, título, subtítulo e a tabela de identificação. Lê os
% campos `\doc...`; linha de campo vazio não sai.
\newcommand{\pb@linha}[2]{\ifx#2\@empty\else #1 & #2 \\\fi}
\newcommand{\rosto}{%
  \thispagestyle{rosto}%
  \begin{adjustwidth}{-\margemmanual}{0pt}
    \raggedright
    \marcadoc[45mm]\par\vspace{16mm}
    {\fontsize{24}{29}\selectfont\bfseries\doctitulo\par}
    \ifx\docsubtitulo\@empty\else
      \vspace{3mm}{\Large\color{cinza}\docsubtitulo\par}\fi
    \vspace{9mm}{\hrule height 0.8pt}\vspace{4mm}
    {\small\renewcommand{\arraystretch}{1.25}%
     \begin{tabular}{@{}>{\bfseries}p{28mm}p{110mm}@{}}
       Documento & \texttt{\doccodigo} \\
       \pb@linha{Versão}{\docversao}%
       Data & \docdata \\
       \pb@linha{Responsável}{\docresponsavel}%
     \end{tabular}\par}
  \end{adjustwidth}
  \vspace{10mm}}

% Histórico de revisões: uma linha por versão, a mais recente em cima.
%   \begin{historico}
%     0.2 & 2026-10-03 & O que mudou \\
%   \end{historico}
\newenvironment{historico}
  {\par\medskip\noindent\small
   \begin{tabular}{@{}p{16mm}p{24mm}p{\dimexpr\linewidth-40mm-4\tabcolsep\relax}@{}}
   \toprule \textbf{Versão} & \textbf{Data} & \textbf{Alteração} \\ \midrule}
  {\bottomrule\end{tabular}\par\medskip}

% ---------------------------------------------------------------------
%  Avisos (ANSI Z535.6) e nota
% ---------------------------------------------------------------------
% Quatro palavras de sinal, em hierarquia fixa. As três primeiras são de
% dano a pessoa; AVISO é de dano a equipamento ou a dado. Cada mensagem
% diz o risco, como evitar e a consequência. NOTA não é aviso: é
% informação de apoio.
\definecolor{sinalperigo}{HTML}{C8102E}
\definecolor{sinalatencao}{HTML}{E87722}
\definecolor{sinalcuidado}{HTML}{F2C200}
\definecolor{sinalaviso}{HTML}{00539B}
\definecolor{sinalnota}{HTML}{5F5F5F}
% #1 palavra, #2 cor da faixa, #3 cor da letra sobre a faixa
\newtcolorbox{pb@mensagem}[3]{enhanced,breakable,sharp corners,
  boxrule=0pt,frame hidden,borderline west={3pt}{0pt}{#2},
  colback=#2!7,left=9pt,right=7pt,top=5pt,bottom=5pt,
  before skip=0.9em,after skip=0.9em,fontupper=\small,
  before upper={\colorbox{#2}{\textcolor{#3}{\footnotesize\bfseries #1}}\enspace}}
\newenvironment{perigo}{\csname pb@mensagem\endcsname{PERIGO}{sinalperigo}{white}}{\csname endpb@mensagem\endcsname}
\newenvironment{atencao}{\csname pb@mensagem\endcsname{ATENÇÃO}{sinalatencao}{black}}{\csname endpb@mensagem\endcsname}
\newenvironment{cuidado}{\csname pb@mensagem\endcsname{CUIDADO}{sinalcuidado}{black}}{\csname endpb@mensagem\endcsname}
\newenvironment{aviso}{\csname pb@mensagem\endcsname{AVISO}{sinalaviso}{white}}{\csname endpb@mensagem\endcsname}
\newenvironment{nota}{\csname pb@mensagem\endcsname{NOTA}{sinalnota}{white}}{\csname endpb@mensagem\endcsname}

% ---------------------------------------------------------------------
%  Procedimento
% ---------------------------------------------------------------------
% Passos numerados, uma ação por passo, no imperativo. A pré-condição vem
% antes da lista e o resultado esperado depois.
\newlist{passos}{enumerate}{2}
\setlist[passos,1]{label=\textbf{\arabic*.},leftmargin=2.2em,itemsep=0.45em,topsep=0.6em}
\setlist[passos,2]{label=\alph*),leftmargin=1.8em,itemsep=0.2em,topsep=0.2em}
\newcommand{\precondicao}[1]{\par\smallskip\noindent\textbf{Antes de começar.}\enspace#1\par}
\newcommand{\resultado}[1]{\par\smallskip\noindent\textbf{Resultado.}\enspace#1\par}

% Bloco na largura útil inteira, margem de título incluída: para tabela
% ou figura que não cabe nos 152 mm do texto.
\newenvironment{larga}{\begin{adjustwidth}{-\margemmanual}{0pt}}{\end{adjustwidth}}
% O mesmo dentro de `table` ou `figure`, onde `adjustwidth` não alcança:
% uma caixa na largura útil, puxada para a margem. Não quebra página.
% Dentro dela a largura é `\linewidth`, e não `\textwidth`.
\newenvironment{tabelalarga}
  {\noindent\hspace*{-\margemmanual}%
   \begin{minipage}{\dimexpr\linewidth+\margemmanual\relax}\centering}
  {\end{minipage}}

% ---------------------------------------------------------------------
%  Campos e figuras
% ---------------------------------------------------------------------
\definecolor{cinza}{gray}{0.40}

% \figuradoc[largura]{caminho}; sem o arquivo, uma caixa marca a pendência
\newcommand{\figuradoc}[2][\linewidth]{%
  \IfFileExists{#2}%
    {\includegraphics[width=#1,height=0.80\textheight,keepaspectratio]{#2}}%
    {\fbox{\parbox[c][45mm][c]{\dimexpr#1-2\fboxsep-2\fboxrule\relax}%
      {\centering\small\color{cinza}\itshape[figura em preparação]}}}}

\pbrestauraarroba

\renewcommand{\doctitulo}{Machine learning and deep learning in the energy sector}
\renewcommand{\docsubtitulo}{Applications, technical foundations and product view}
\renewcommand{\docresponsavel}{Mapeamento Energético}
\renewcommand{\docdata}{2026-09}

% English labels for the cover block, the footer and the note box.
\makeatletter
\renewcommand{\pb@identificacao}{%
  \texttt{\doccodigo}%
  \ifx\docversao\@empty\else\enspace\textperiodcentered\enspace version \docversao\fi
  \enspace\textperiodcentered\enspace\docdata}
\fancyfoot[R]{\footnotesize\thepage\ of \pageref*{LastPage}}
\renewcommand{\rosto}{%
  \thispagestyle{rosto}%
  \begin{adjustwidth}{-\margemmanual}{0pt}
    \raggedright
    \marcadoc[45mm]\par\vspace{16mm}
    {\fontsize{24}{29}\selectfont\bfseries\doctitulo\par}
    \ifx\docsubtitulo\@empty\else
      \vspace{3mm}{\Large\color{cinza}\docsubtitulo\par}\fi
    \vspace{9mm}{\hrule height 0.8pt}\vspace{4mm}
    {\small\renewcommand{\arraystretch}{1.25}%
     \begin{tabular}{@{}>{\bfseries}p{28mm}p{110mm}@{}}
       Document & \texttt{\doccodigo} \\
       \pb@linha{Version}{\docversao}%
       Date & \docdata \\
       \pb@linha{Owner}{\docresponsavel}%
     \end{tabular}\par}
  \end{adjustwidth}
  \vspace{10mm}}
\renewenvironment{nota}{\csname pb@mensagem\endcsname{NOTE}{sinalnota}{white}}{\csname endpb@mensagem\endcsname}
\makeatother

% en-US numbers: decimal point, comma as thousands separator.
\sisetup{
  output-decimal-marker = {.},
  group-separator       = {,},
  range-phrase          = {\,to\,},
  list-pair-separator   = {\ and\ },
  list-final-separator  = {\ and\ },
}

% `L` here is a left-aligned `X` column, and the tables of this document use
% it that way; it overrides the `L{width}` of the base preamble.
\newcolumntype{L}{>{\raggedright\arraybackslash}X}

\definecolor{notecolor}{RGB}{110,110,110}

% Product view: subsection with the full title in the table of contents and
% the short form in the header, where the full title would run into the
% document title.
\newcommand{\visaoproduto}[1]{%
  \let\subsectionmarkbase\subsectionmark
  \renewcommand{\subsectionmark}[1]{\markright{\thesubsection\ Product view}}%
  \subsection{Product view: #1}%
  \let\subsectionmark\subsectionmarkbase}

\begin{document}

\rosto

\begin{abstract}
This document describes four application areas of machine learning (ML) and
deep learning (DL) in the power sector: (1)~forecasting of load, renewable
generation and price; (2)~grid operation and optimization; (3)~fault
detection, protection and predictive maintenance; and (4)~distribution and the
end consumer. For each area it presents the business problem, the
mathematical formulation, the model families, the required data, the
evaluation metrics and the practical challenges. Each section ends with a
product view: who the user is, which problem is addressed, the minimum viable
scope, the reference architecture, the success metrics, the business model
and the risks. The Brazilian regulatory and operational context (ONS, CCEE,
ANEEL, distribution utilities) is used as the reference where applicable.
\end{abstract}

\tableofcontents
\newpage

% ===========================================================================
\section{Introduction}
% ===========================================================================

The power sector is undergoing three simultaneous transformations that widen
the scope for ML/DL:

\begin{enumerate}
  \item \textbf{Decarbonization}: the growing share of wind and solar
    sources, which are variable and non-dispatchable, shifts the central
    problem of operation from ``serving a predictable load with controllable
    generation'' to ``balancing uncertainties on both sides''.
  \item \textbf{Decentralization}: distributed micro and mini generation
    (MMGD, from the Portuguese \textit{micro e minigeração distribuída}),
    batteries and electric vehicles turn the consumer into an active agent
    and make the distribution grid bidirectional.
  \item \textbf{Digitalization}: smart meters, phasor measurement units
    (PMUs), sensors on assets, drones and satellite imagery generate volumes
    of data that classical methods do not exploit well.
\end{enumerate}

Despite the large volume of research, recent reviews converge on one
diagnosis: there are many pilots and few systems at scale, and the main
barrier is not model accuracy but integration with legacy systems,
reliability and acceptance by operators~\cite{energyinformatics2025,trustworthy2024}.
For this reason, each section of this document treats the model and the
product as equally important problems.

% ===========================================================================
\section{Forecasting: load, renewable generation and price}
% ===========================================================================

\subsection{Context and importance}

Forecasting is one of the most mature applications of ML in the sector.
Operational and commercial decisions depend on an estimate of the future:

\begin{itemize}
  \item the \textbf{system operator} (in Brazil, ONS) needs to forecast load
    and renewable generation to schedule dispatch, size reserves and assess
    security;
  \item \textbf{generators and trading companies} need to forecast production
    and price to contract, hedge and avoid exposure to the short-term
    market;
  \item \textbf{distribution utilities} need to forecast demand for energy
    purchases in auctions, expansion planning and management of transformer
    loading.
\end{itemize}

A frequently cited case of value generated is that of DeepMind with Google: a
neural network trained on weather forecasts and turbine history started to
forecast the generation of 700\,MW of wind farms 36 hours ahead. With that
forecast, Google started to offer the energy in the day-ahead market, and the
company reports an increase of about 20\% in the value of the wind
energy~\cite{deepmindwind}.

\subsection{Forecast horizons}

\begin{table}[h]
\centering
\caption{Forecast horizons and their uses.}
\begin{tabularx}{\textwidth}{l l L}
\toprule
\textbf{Horizon} & \textbf{Typical range} & \textbf{Main use} \\
\midrule
Very short term & seconds to 1\,h & frequency control, ramps, intra-hour
  dispatch, inverter control \\
Short term & 1\,h to 7 days & daily operation scheduling, day-ahead market,
  unit commitment \\
Medium term & weeks to months & operation planning (PMO), maintenance,
  reservoir management, contracting \\
Long term & years & generation and transmission expansion, auctions, tariff
  planning \\
\bottomrule
\end{tabularx}
\end{table}

Each horizon has different data, models and metrics. At very short horizons,
local high-resolution data predominate (measurements, sky cameras, satellite
images). In the short term, numerical weather prediction (NWP) is the main
explanatory variable. In the medium and long term, macroeconomic and climate
scenarios enter, and the problem is closer to structural modeling than to
pure ML.

\subsection{Problem formulation}

Let $y_t$ be the target variable (load, generated power or price) at time
$t$, $\mathbf{x}_t$ a vector of observed covariates (temperature, measured
irradiance, etc.) and $\mathbf{z}_{t+h}$ covariates known or forecast for the
future (calendar, holidays, NWP). The point forecast with horizon $h$ and
context window $L$ is
\begin{equation}
  \hat{y}_{t+h} = f_\theta\big(y_{t-L+1:t},\; \mathbf{x}_{t-L+1:t},\;
  \mathbf{z}_{t+1:t+h}\big).
\end{equation}

In practice, the decision rarely depends on the expected value alone. Sizing
reserve, bidding in a market or deciding whether to start a thermal plant
requires knowing the \emph{distribution} of the future. For this reason the
trend is toward \textbf{probabilistic forecasting}, usually through quantiles
$q_\tau$ trained with the pinball loss:
\begin{equation}
  \mathcal{L}_\tau(y, q_\tau) =
  \max\big(\tau\,(y - q_\tau),\; (\tau - 1)\,(y - q_\tau)\big),
  \qquad \tau \in (0,1).
\end{equation}
The quality of the full distribution is measured by the Continuous Ranked
Probability Score (CRPS), where $F$ is the forecast distribution:
\begin{equation}
  \mathrm{CRPS}(F, y) = \int_{-\infty}^{\infty}
  \big(F(z) - \mathbb{1}\{z \ge y\}\big)^2 \, dz.
\end{equation}
The two measures are related: the CRPS is the pinball loss integrated over
all quantiles,
\begin{equation}
  \mathrm{CRPS}(F, y) = 2 \int_0^1 \mathcal{L}_\tau\big(y, F^{-1}(\tau)\big)\, d\tau,
\end{equation}
and it reduces to the absolute error when the forecast is a point forecast.

The quantile that matters depends on the decision. If each MWh
under-forecast costs $c_{-}$ and each MWh over-forecast costs $c_{+}$, the
forecast that minimizes the expected cost is the quantile
\begin{equation}
  \tau^{*} = \frac{c_{-}}{c_{-} + c_{+}},
\end{equation}
which coincides with the median only when the two costs are equal. A fixed
set of quantiles (P10, P50, P90) serves the decision only if $\tau^{*}$ is
one of them.

\subsection{Load forecasting}

\paragraph{Load drivers.} Aggregate electrical load has strong and relatively
stable patterns:
\begin{itemize}
  \item \textbf{multiple seasonality}: daily, weekly and annual;
  \item \textbf{calendar effects}: holidays, bridge days, daylight saving
    time (where it exists), large events;
  \item \textbf{weather}: temperature is the main variable (a nonlinear,
    U-shaped relationship, because of air conditioning and heating),
    followed by humidity, apparent temperature and cloud cover;
  \item \textbf{economy}: industrial and commercial activity, shocks such as
    the 2020 pandemic.
\end{itemize}

\paragraph{Gross load versus net load.} With the expansion of solar MMGD,
what the operator observes at the connection points is the \emph{net load}:
consumption minus behind-the-meter generation. On sunny days, the net load
drops at midday and rises abruptly in the late afternoon, forming the
so-called ``duck curve''. Forecasting the net load directly mixes two
phenomena with different dynamics. For this reason, the common practice is
to decompose:
\begin{equation}
  \text{net load}_t = \text{gross load}_t - \text{MMGD}_t,
\end{equation}
forecasting each term with its own covariates. For this reason, ONS maintains
a dedicated system for forecasting photovoltaic MMGD generation, with a
72-hour horizon~\cite{onssolar}.

The decomposition has two conditions. The first is one of identification:
only the net load is measured. Behind-the-meter MMGD has to be estimated
(registered installed capacity and irradiance), and the gross load is
reconstructed from that estimate, so the error of the estimate enters both
terms. The second is one of variance: with $e_b$ and $e_m$ the forecast
errors of the gross load and of MMGD,
\begin{equation}
  \mathrm{Var}(e_b - e_m) = \mathrm{Var}(e_b) + \mathrm{Var}(e_m)
  - 2\,\mathrm{Cov}(e_b, e_m),
\end{equation}
and the decomposition outperforms the direct forecast only if this variance
is smaller than the error variance of the direct model. The comparison is
made by backtesting.

\paragraph{Hierarchical forecasting.} Load exists at several levels
(subsystem, area, substation, feeder, transformer, consumer). Independent
forecasts at each level do not add up coherently. \textbf{Hierarchical
reconciliation} techniques (such as MinT) adjust the forecasts so that the
levels are consistent with each other, and they tend to improve accuracy at
all levels.

\subsection{Wind generation forecasting}

The power extracted from the wind grows with the cube of the speed:
\begin{equation}
  P = \tfrac{1}{2}\,\rho\, A\, C_p(\lambda, \beta)\, v^3,
\end{equation}
where $\rho$ is the air density, $A$ the swept area, $C_p$ the power
coefficient and $v$ the wind speed. With constant $C_p$, the relative error
triples: $\delta P / P \approx 3\,\delta v / v$, and an error of 10\% in the
wind corresponds to 33\% in power ($1.1^3 = 1.331$). The relationship holds
in the cubic region of the power curve; above rated speed the power is
limited and the sensitivity to wind falls to zero.

The cube also biases the conversion. For wind with mean $\bar{v}$ and
turbulence intensity $I = \sigma_v / \bar{v}$,
\begin{equation}
  \mathbb{E}\big[v^3\big] \approx \bar{v}^3\,\big(1 + 3I^2\big),
\end{equation}
so applying the power curve to the mean wind forecast underestimates the mean
power: with $I = 0.15$, $\mathbb{E}[v^3]$ exceeds $\bar{v}^3$ by 6.75\%. Near
rated speed the curve is concave and the bias changes sign.

The typical architecture of short-term forecasting has two stages:
\begin{enumerate}
  \item \textbf{Weather forecast}: global or regional NWP models provide wind
    at several heights, direction, temperature and pressure. AI-based weather
    models (such as GraphCast, GenCast and Aurora) have come to compete with
    physical NWP on several metrics, at a much lower computational cost.
  \item \textbf{Conversion model}: an ML model (usually gradient boosting or
    neural networks) learns the relationship between the forecast weather
    and the measured power of the farm. This model corrects NWP biases, wake
    effects, unavailability and losses.
\end{enumerate}

ONS developed its own wind forecasting model, aimed at the Northeast and
South regions, to reduce the need for power reserve~\cite{onseolico}.

\begin{nota}
\textbf{Caution with curtailment.} When the operator restricts the generation
of a farm (in Brazil, the so-called \textit{constrained-off}, which has
become relevant for wind and solar plants in the Northeast), the measured
power no longer reflects the available resource. A model trained on these
data without filtering them underestimates generation. The periods of
restriction need to be flagged and excluded, or the \emph{available} power
modeled instead of the \emph{delivered} power.
\end{nota}

\subsection{Solar generation forecasting}

The irradiance that reaches the surface is dominated by a deterministic
component (the position of the sun) and a stochastic component (clouds and
aerosols). The standard normalization is the \textbf{clear-sky index}:
\begin{equation}
  k^{\mathrm{cs}}_t = \frac{\mathrm{GHI}_t}{\mathrm{GHI}^{\mathrm{cs}}_t},
\end{equation}
where $\mathrm{GHI}^{\mathrm{cs}}$ is the global horizontal irradiance
computed by a clear-sky model. Forecasting $k^{\mathrm{cs}}_t$ instead of GHI
removes the astronomical seasonality and restricts the problem to the
stochastic component: the clouds. The clear-sky index is not the clearness
index, which divides GHI by the extraterrestrial irradiance and is usually
written $k_t$.

The denominator tends to zero at sunrise and sunset, and the ratio amplifies
the measurement error at those times. For this reason the index is computed
and evaluated only above a minimum solar elevation (of the order of
$5^\circ$). The index also exceeds 1 in episodes of reflection at cloud
edges.

The data sources and models vary with the horizon:
\begin{itemize}
  \item \textbf{minutes}: all-sky imagers and CNNs that estimate cloud
    motion;
  \item \textbf{up to 6 hours}: geostationary satellite images (in Brazil,
    GOES-East) and video/spatiotemporal models (ConvLSTM, U-Net);
  \item \textbf{day ahead and beyond}: NWP as input to ML models.
\end{itemize}

The conversion from irradiance to power involves module temperature
(efficiency falls with heat), orientation, shading, soiling and the inverter
power limit (clipping). A Brazilian study uses LSTM for hourly forecasting of
photovoltaic generation~\cite{pvlstm}.

\subsection{Price forecasting}

In bid-based markets (Europe, USA), deep networks are among the methods with
the lowest reported error for day-ahead price forecasting. Brazil is
different: the Settlement Price for Differences (PLD, \textit{Preço de
Liquidação das Diferenças}) does not result from bids; it is computed by CCEE
from the chain of hydrothermal optimization models (NEWAVE, DECOMP and
DESSEM, developed by CEPEL). Since 2021 the PLD has been hourly and computed
by DESSEM.

This opens two ML strategies:
\begin{enumerate}
  \item \textbf{Direct PLD forecasting} from variables such as inflows,
    storage, load and forecast renewable generation;
  \item \textbf{Emulation of the official models} (surrogate modeling):
    training a network to approximate the output of DESSEM/DECOMP as a
    function of the inputs, which allows thousands of scenarios to be run in
    seconds. It is useful for the risk analysis of trading companies.
\end{enumerate}

\subsection{Model families}

\begin{table}[h]
\small
\caption{Model families for forecasting in energy.}
\begin{tabelalarga}
\begin{tabularx}{\linewidth}{l L L}
\toprule
\textbf{Family} & \textbf{Strengths} & \textbf{Limitations} \\
\midrule
Statistical (ARIMA, ETS, regression) & interpretable, low cost, suitable as
  baselines & capture nonlinearities and many covariates poorly \\
Gradient boosting (XGBoost, LightGBM) & good performance on tabular data
  with rich covariates; widely used in industry & requires feature
  engineering (lags, windows) \\
RNN (LSTM, GRU) & model temporal dependence directly & slower training; do
  not consistently outperform GBM \\
Temporal CNNs (TCN) & parallelizable, long receptive field & less used with
  future covariates \\
N-BEATS, N-HiTS & deep MLPs with interpretable decomposition & univariate in
  the original form \\
Transformers (TFT, PatchTST) & heterogeneous covariates, multiple horizons,
  quantiles (TFT)~\cite{tft} & require more data and tuning \\
Foundation models (Chronos, TimesFM, Moirai) & zero-shot forecasting without
  local training~\cite{chronos,timesfm} & limited covariate support;
  inference cost \\
\bottomrule
\end{tabularx}
\end{tabelalarga}
\end{table}

\paragraph{Time series foundation models.} These are models pretrained on
large collections of heterogeneous time series, able to forecast unseen
series without retraining. The line started with TimeGPT (2023) and continued
with Lag-Llama, Chronos, TimesFM and, more recently, Sundial and
Kairos~\cite{fmreview}. Results from 2025--2026 in energy:
\begin{itemize}
  \item in load benchmarks at several grid levels, time series Transformers
    outperformed the established methods~\cite{loadbenchmark};
  \item foundation models run on consumer hardware with good zero-shot
    performance, but with sensitivity to context length and to distribution
    shift~\cite{tsfmconsumer};
  \item the incorporation of covariates (weather, calendar) improves the
    performance of the models and allows the forecasts to be
    explained~\cite{tsfmxai};
  \item at low voltage, where the series are noisy and short, the models
    were applied to probabilistic peak forecasting~\cite{lvpeak};
  \item energy-specific models, such as GridFM, jointly forecast load,
    price, emissions and renewables, incorporating physical
    constraints~\cite{gridfm}.
\end{itemize}
The main practical benefit is the \textbf{cold start}: new farms, new
feeders or new customers receive forecasts on the first day, before there is
history to train a dedicated model.

\subsection{Evaluation}

\paragraph{Metrics.}
\begin{itemize}
  \item \textbf{MAE, RMSE}: in the unit of the variable; the RMSE penalizes
    large errors more, which matters for ramps.
  \item \textbf{MAPE}: popular for load, but \emph{unsuitable for solar},
    which is zero at night (division by zero) and has small values at the
    start and end of the day. For renewables, the error normalized by
    installed capacity is used (nMAE, nRMSE).
  \item \textbf{Skill score}: relative gain over a naive reference
    (persistence or climatology),
    $\mathrm{SS} = 1 - \mathrm{RMSE}_{\text{model}}/\mathrm{RMSE}_{\text{ref}}$.
    Without this comparison, absolute error numbers are not very
    informative.
  \item \textbf{Probabilistic}: pinball loss, CRPS, coverage and width of
    the intervals (calibration and sharpness).
  \item \textbf{Value metrics}: avoided imbalance cost, saved reserve,
    additional revenue. These are the metrics that matter to the business.
\end{itemize}

\paragraph{Validation.} Two common errors artificially inflate the results:
\begin{enumerate}
  \item \textbf{Random cross-validation} on time series leaks the future
    into training. The appropriate procedure is backtesting with a rolling
    origin, which respects temporal order.
  \item \textbf{Covariate leakage}: using the \emph{observed} weather in
    place of the \emph{forecast} available at decision time. In production
    only the forecast exists, issued at a specific time. Training needs to
    reproduce what would have been available at each moment.
\end{enumerate}

\subsection{Challenges}
\begin{itemize}
  \item \textbf{Extreme events}: heat waves, storms and blackouts are rare
    in the history and are the ones with the largest consequences.
  \item \textbf{Structural changes}: growth of MMGD, electrification, changes
    in habits. Models need drift monitoring and continuous retraining.
  \item \textbf{Data quality}: missing measurements, telemetry errors,
    unflagged curtailment.
  \item \textbf{User trust}: the operator needs to understand why the
    forecast changed. Explainability (covariate importance, decomposition)
    is a requirement.
\end{itemize}

\visaoproduto{Forecasting platform for renewables and load}

\paragraph{Problem.}
Wind and solar generators are exposed to imbalance penalties, curtailment and
the short-term price. Trading companies need reliable forecasts to contract.
Distribution utilities and large consumers have an increasingly uncertain
net load because of MMGD. Many of these agents use spreadsheets, generic
forecasts from foreign vendors or persistence.

\paragraph{Users and personas.}
\begin{itemize}
  \item \textbf{Commercial operations analyst of a generator}: needs the
    generation forecast per farm for scheduling and for declarations to the
    operator.
  \item \textbf{Trader at a trading company}: needs forecasts of aggregate
    generation, load and PLD, with uncertainty, to decide positions.
  \item \textbf{Distribution utility planner}: needs the load per
    substation/feeder and the MMGD per region.
\end{itemize}

\paragraph{Value proposition.}
Probabilistic forecasts per asset, with horizons from minutes to 15 days,
available through an API and a dashboard, calibrated for the Brazilian
context (NWP and satellite over the territory, national calendar, curtailment
flagging), with a measured gain over the reference that the customer already
uses.

\paragraph{MVP scope.}
\begin{enumerate}
  \item Day-ahead and intraday forecasting of wind and solar generation for
    farms with available telemetry.
  \item Quantiles (P10, P50, P90), not only a point value.
  \item Dashboard with a daily comparison of forecast versus actual and
    skill score against persistence.
  \item REST API for integration with the customer's systems.
  \item Onboarding of new farms with a zero-shot foundation model, replaced
    by a fitted model once there is history.
\end{enumerate}

\paragraph{Evolution (roadmap).}
\begin{itemize}
  \item \textbf{V2}: nowcasting with satellite images (0--6\,h); automatic
    detection of curtailment and unavailability.
  \item \textbf{V3}: forecasting of PLD and of load per submarket; joint
    scenarios (generation, load, price) for risk analysis.
  \item \textbf{V4}: optimization of bids and of battery dispatch coupled to
    the forecasts (from forecast to decision).
\end{itemize}

\paragraph{Reference architecture.}
\begin{itemize}
  \item \textbf{Ingestion}: NWP (multiple providers), satellite, SCADA
    telemetry of the farms, open data from ONS and CCEE.
  \item \textbf{Feature store} with temporal versioning (``what was known at
    each instant'') to avoid leakage.
  \item \textbf{Training} per asset or global (one model for several farms),
    with automatic selection by backtesting.
  \item \textbf{Inference service} scheduled according to the arrival of
    each NWP run.
  \item \textbf{Monitoring} of error in production, drift and alerts for
    missing data.
\end{itemize}

\paragraph{Success metrics.}
\begin{itemize}
  \item Technical: nMAE and CRPS per asset; skill score over persistence and
    over the customer's current vendor.
  \item Business: reduction in imbalance cost (R\$/MWh), additional revenue,
    analyst hours saved.
  \item Product: onboarding time for a new farm, API availability, customer
    retention.
\end{itemize}

\paragraph{Business model.}
Subscription per monitored MW or per asset, with tiers (day ahead; intraday;
nowcasting; price). A free or low-cost pilot in shadow mode, compared with
the customer's current forecast, is the main sales instrument: the gain
appears in the customer's own numbers.

\paragraph{Risks.}
\begin{itemize}
  \item Dependence on third-party NWP (cost, availability, version changes
    that alter the input distribution).
  \item Access to the customer's telemetry (bureaucracy, information
    security).
  \item A market with established international competitors; differentiation
    needs to come from local accuracy and integration with the Brazilian
    regulatory context.
\end{itemize}

% ===========================================================================
\section{Grid operation and optimization}
% ===========================================================================

\subsection{Context}

Operating a power system in real time involves repeatedly solving analysis
and optimization problems:
\begin{itemize}
  \item \textbf{Power flow}: given the state of generation and load, compute
    voltages and flows across the whole grid;
  \item \textbf{Contingency analysis}: check whether the system remains
    secure after the loss of any element ($N-1$ criterion), which requires
    thousands of power flows;
  \item \textbf{Optimal power flow (OPF)}: choose the lowest-cost dispatch
    that respects all physical limits;
  \item \textbf{State estimation}: reconstruct the most likely state of the
    grid from noisy and incomplete measurements;
  \item \textbf{Control}: voltage and reactive power, redispatch, switching
    (topology), batteries.
\end{itemize}

With more renewables, these problems need to be solved more often, for more
scenarios and with more uncertainty. Classical solvers are reliable, but the
computational cost grows rapidly with the size of the grid and the number of
scenarios. ML enters here mainly as an \textbf{accelerator} and as
\textbf{decision support}, not as a substitute for the physical methods.

\subsection{Formulation: AC power flow}

For each bus $i$ of the grid, the active and reactive power injections are
\begin{align}
  P_i &= \sum_{j} |V_i||V_j| \big(G_{ij}\cos\theta_{ij} + B_{ij}\sin\theta_{ij}\big), \\
  Q_i &= \sum_{j} |V_i||V_j| \big(G_{ij}\sin\theta_{ij} - B_{ij}\cos\theta_{ij}\big),
\end{align}
where $V_i$ is the voltage phasor, $\theta_{ij} = \theta_i - \theta_j$ and
$G_{ij} + jB_{ij}$ are the elements of the nodal admittance matrix. The
system is nonlinear and is solved by Newton--Raphson.

The OPF adds an objective function and constraints:
\begin{equation}
\begin{aligned}
  \min_{p, q, |V|, \theta} \quad & \sum_{g} c_g(p_g) \\
  \text{s.t.} \quad
  & P_i = \sum_{g \in \mathcal{G}_i} p_g - P^{d}_i, \qquad
    Q_i = \sum_{g \in \mathcal{G}_i} q_g - Q^{d}_i, \\
  & V_i^{\min} \le |V_i| \le V_i^{\max}, \\
  & p_g^{\min} \le p_g \le p_g^{\max}, \qquad
    q_g^{\min} \le q_g \le q_g^{\max}, \\
  & |S_{ij}| \le S_{ij}^{\max}, \qquad |S_{ji}| \le S_{ij}^{\max}, \\
  & \theta_{\mathrm{ref}} = 0.
\end{aligned}
\end{equation}
Here $p_g$ and $q_g$ are the active and reactive power of generator $g$,
$\mathcal{G}_i$ is the set of generators at bus $i$, $P^{d}_i$ and $Q^{d}_i$
are the load at the bus, and $P_i$ and $Q_i$ are the injections of the power
flow equations. The flow limit applies at both terminals of the line,
because losses make $|S_{ij}| \ne |S_{ji}|$, and the reference bus fixes the
origin of the angles. The problem is nonconvex and, in large systems with
security constraints, costly to solve within real-time deadlines.

\subsection{Graph neural networks (GNNs)}

A power grid can be represented as a graph: buses are nodes, lines and
transformers are edges. GNNs process this type of data by message passing:
\begin{equation}
  \mathbf{h}_i^{(k+1)} = \phi\Big(\mathbf{h}_i^{(k)},\;
  \bigoplus_{j \in \mathcal{N}(i)} \psi\big(\mathbf{h}_i^{(k)}, \mathbf{h}_j^{(k)}, \mathbf{e}_{ij}\big)\Big),
\end{equation}
where $\mathbf{h}_i$ is the representation of bus $i$ (load, generation,
type), $\mathbf{e}_{ij}$ are attributes of the line (impedance, limit) and
$\bigoplus$ is a permutation-invariant aggregation (sum, mean). After $K$
layers, the representation of each bus incorporates its $K$-hop
neighborhood.

This structure brings two advantages over MLPs that receive a fixed vector:
\begin{enumerate}
  \item \textbf{Topological generalization}: the same GNN accepts as input
    the grid without a line or with the topology changed, because the
    parameters are shared among nodes and edges.
  \item \textbf{Scalability}: the cost grows with the number of edges, and
    models trained on smaller grids can be applied to larger grids.
\end{enumerate}

Accepting the input does not guarantee a correct output. Power flow is a
global problem: in the linear approximation, the voltages are
$\mathbf{V} = \mathbf{Y}^{-1}\mathbf{I}$, and $\mathbf{Y}^{-1}$ is dense even
when $\mathbf{Y}$ is sparse, so the voltage at one bus depends on the
injections of the whole grid. A GNN with $K$ smaller than the diameter of
the graph does not represent this dependence. Two limits to the advantages
above follow from this: the outage of a line redistributes flow beyond the
$K$ hops, and a larger grid has a larger diameter, which requires more
layers or global aggregation. The error under these two conditions is
measured by testing with topologies and sizes outside the training set.

\paragraph{Main applications.}
\begin{itemize}
  \item \textbf{Power flow surrogate}: predicting voltages and flows
    directly, in milliseconds, for fast contingency screening. A recent
    study reports GNNs that scale to larger systems with a computation time
    lower than that of traditional solvers~\cite{gnnpf}.
  \item \textbf{OPF warm start}: the GNN provides an initial point close to
    the solution, and the classical solver converges in fewer iterations,
    which preserves the feasibility guarantee~\cite{gnnopfinit}.
  \item \textbf{Network reconfiguration}: predicting the topology before the
    optimization stage, which reduces the combinatorial search
    space~\cite{gnnreconf}.
\end{itemize}

\subsection{Learning to optimize and the feasibility problem}

Training a network to map inputs (load, costs) directly to OPF solutions
(optimization proxies) has a limitation: the network does not guarantee that
the solution respects the constraints. A violation of a thermal or voltage
limit is not a ``small error''; it is an unacceptable solution. The main
strategies are:
\begin{itemize}
  \item \textbf{Penalization} of the violations in the loss function (low
    implementation cost, no guarantee);
  \item \textbf{Lagrangian duality} during training, adjusting the weights of
    the constraints;
  \item \textbf{Correction and completion}: the network predicts part of the
    variables and the others are obtained by solving the equalities,
    followed by correction steps for the inequalities (as in
    DC3~\cite{dc3});
  \item \textbf{Projection} of the output onto the feasible set, or use of
    the prediction only as a warm start for an exact solver.
\end{itemize}
In industrial practice, the most accepted approach is the last one: ML
accelerates, the solver certifies.

\subsection{Reinforcement learning (RL)}

Sequential control problems, such as deciding switching actions, redispatch
or battery cycles over time, can be formulated as Markov decision processes.
The agent observes the state $s_t$ (flows, voltages, forecasts), chooses an
action $a_t$ and receives a reward $r_t$ (avoided cost, security margin). The
objective is
\begin{equation}
  \max_\pi \; \mathbb{E}_\pi\Big[\sum_{t=0}^{T} \gamma^t r(s_t, a_t)\Big].
\end{equation}

\paragraph{Topology control.} The \textit{Learning to Run a Power Network}
(L2RPN) competition, organized by RTE (the French operator), popularized the
use of RL to keep the grid operating through topology changes in
substations, redispatch and curtailment~\cite{l2rpn}. Agents that combine
GNNs with RL (graph RL), sometimes pretrained by imitation learning and then
refined with PPO, were proposed to handle topology changes caused by extreme
weather or maintenance~\cite{graphrlsurvey,ppognn}.

\paragraph{Batteries and microgrids.} RL is used to decide when to charge
and discharge batteries considering price, generation forecast and
degradation. PINN-based surrogates allow agents to be trained without
interacting with the real system~\cite{rlpinn}.

\paragraph{Safe RL.} In physical systems, exploring bad actions has a real
cost. The standard formulation is the constrained MDP:
\begin{equation}
  \max_\pi \; \mathbb{E}_\pi\Big[\sum_t \gamma^t r_t\Big]
  \quad \text{s.t.} \quad
  \mathbb{E}_\pi\Big[\sum_t \gamma^t c_t\Big] \le d,
\end{equation}
where $c_t$ measures violations (overload, undervoltage). The constraint
holds in expectation and with discounting: a policy satisfies it even with
occasional overloads, as long as the discounted mean stays below $d$. A
physical limit calls for a per-instant constraint, in chance form,
\begin{equation}
  \Pr\big(c_t > 0\big) \le \epsilon \quad \text{for all } t,
\end{equation}
or a constraint on the tail of the cost distribution (CVaR). Other
techniques include shielding (a module checks and blocks dangerous actions
before execution) and offline training from historical operator
data~\cite{saferl}.

\subsection{Physics-informed neural networks (PINNs)}

PINNs incorporate physical equations into the loss function:
\begin{equation}
  \mathcal{L} = \mathcal{L}_{\text{data}} + \lambda\, \mathcal{L}_{\text{physics}},
\end{equation}
where $\mathcal{L}_{\text{physics}}$ penalizes the residual of the equations
(for example, of the power flow equations or of the differential equations
of machine dynamics)~\cite{pinn}. The reported benefits are a lower need for
labeled data, closer adherence to physical laws and better extrapolation.
PINNs are used in state estimation, transient stability analysis and as
simulators for training RL agents.

\subsection{Challenges}
\begin{itemize}
  \item \textbf{Certification and liability}: operators and regulators
    require guarantees; a model that is correct in 99.9\% of the cases can
    still fail in the critical contingency. In a sweep of \num{5000}
    contingencies with independent errors, the probability of at least one
    error is $1 - 0.999^{5000} \approx 99.3\%$.
  \item \textbf{Generalization}: the real grid changes (new lines, new
    plants), and models trained on past scenarios may not cover future
    situations~\cite{gridgeneralization}.
  \item \textbf{Integration with legacy EMS}: energy management systems are
    critical, closed and have long update cycles.
  \item \textbf{Data}: detailed grid models and operational data are
    sensitive; research depends heavily on synthetic test systems (IEEE,
    PGLib).
  \item \textbf{Human factor}: full autonomy is not realistic in the short
    term. The viable role is that of a copilot for the operator.
\end{itemize}

\visaoproduto{Grid operation copilot}

\paragraph{Problem.}
Control center operators (in transmission or distribution) need to make fast
decisions in the face of a growing number of scenarios: renewable
variability, curtailment, contingencies, transmission constraints. The
analysis tools are slow for exploring many alternatives in real time, and
the experience of senior operators is not codified.

\paragraph{Users and personas.}
\begin{itemize}
  \item \textbf{Real-time operator} of a transmission company or of the
    distribution operation center (COD);
  \item \textbf{Electrical studies engineer} who prepares the operation of
    the following day;
  \item \textbf{Operator of renewable farms and batteries} who needs to
    decide dispatch and storage.
\end{itemize}

\paragraph{Value proposition.}
Contingency screening and recommendation of corrective actions in seconds,
with each recommendation verified by a physical solver before it is
presented. The operator remains in command; the system increases the number
of scenarios the operator can evaluate.

\paragraph{MVP scope.}
\begin{enumerate}
  \item GNN-based power flow surrogate, trained on simulations of the
    customer's grid model.
  \item Fast screening of $N-1$ contingencies for the next intervals, using
    load and renewable forecasts, with the critical cases confirmed by a
    full power flow.
  \item ``Hot spot'' dashboard: which lines and buses are at risk and when.
  \item Operation in \textbf{shadow mode}: the system runs in parallel,
    without acting, and the recommendations are recorded against the actual
    decisions.
\end{enumerate}

\paragraph{Evolution (roadmap).}
\begin{itemize}
  \item \textbf{V2}: recommendation of corrective actions (redispatch,
    switching), ranked and verified by the solver.
  \item \textbf{V3}: RL agent trained on historical decisions and in
    simulation, suggesting sequences of actions with security constraints.
  \item \textbf{V4}: coordinated optimization of batteries and curtailment
    in real time for hybrid farms.
\end{itemize}

\paragraph{Reference architecture.}
\begin{itemize}
  \item \textbf{Digital twin} of the grid (electrical model synchronized
    with the current topology).
  \item \textbf{Scenario generator} to create synthetic training data from
    the physical solver.
  \item \textbf{Models}: power flow GNN; contingency ranker; RL policy with
    a shield.
  \item \textbf{Verification layer}: every output passes through the solver
    before reaching the operator.
  \item \textbf{Integration} with SCADA/EMS through standard protocols (such
    as ICCP and CIM), in read-only mode at first.
\end{itemize}

\paragraph{Success metrics.}
\begin{itemize}
  \item Technical: voltage and flow error of the surrogate; recall of
    critical contingencies (missing a critical one is the unacceptable
    error); speed-up relative to the solver.
  \item Operational: reduction of avoidable curtailment, reduction of
    violations, response time to events.
  \item Adoption: percentage of recommendations accepted by the operators.
\end{itemize}

\paragraph{Business model.}
Software license per control center, with a service for deploying and
calibrating the digital twin. For transmission companies and distribution
utilities, R\&D projects regulated by ANEEL are an entry route for pilots.

\paragraph{Risks.}
\begin{itemize}
  \item Long sales cycle and risk aversion of the customers.
  \item Access to the grid model and to operational data.
  \item Cybersecurity requirements for any integration with control systems.
  \item Liability in case of a wrong recommendation; the
    ``human in control'' design and the physical verification mitigate it
    but do not eliminate it.
\end{itemize}

% ===========================================================================
\section{Fault detection, protection and predictive maintenance}
% ===========================================================================

\subsection{Context}

Equipment failures cause supply interruptions, damage to costly assets and
safety risks. In Brazil, ANEEL measures continuity of supply through
collective indicators (DEC and FEC, the equivalent duration and frequency of
interruption per consumer unit) and individual ones (DIC, FIC, DMIC), and
the violation of the individual limits generates compensation paid to
consumers. Reducing the frequency and duration of failures therefore has
direct financial value.

Maintenance evolves in four stages:
\begin{enumerate}
  \item \textbf{Corrective}: repair after failure;
  \item \textbf{Preventive}: periodic interventions by calendar;
  \item \textbf{Condition-based}: intervene when measured indicators exceed
    limits;
  \item \textbf{Predictive}: estimate when the failure is expected to occur
    and act before it.
\end{enumerate}
ML enables the move from the third to the fourth stage at large scale.

\subsection{Fault detection, classification and location}

A fault (short circuit) on a line produces voltage and current transients
recorded by digital relays and fault recorders. The analysis involves three
tasks~\cite{faultsurvey}:
\begin{itemize}
  \item \textbf{Detection}: was there a fault?
  \item \textbf{Classification}: of which type (phase-to-ground,
    phase-to-phase, three-phase, etc.)?
  \item \textbf{Location}: at what distance from the terminal?
\end{itemize}

\paragraph{Approaches.}
\begin{itemize}
  \item \textbf{Feature extraction + classifier}: wavelet or Fourier
    transform of the signals, followed by SVM, trees or shallow networks.
  \item \textbf{End-to-end learning}: 1D CNNs on the raw waveforms or 2D CNNs
    on time-frequency representations (spectrograms, scalograms).
  \item \textbf{PMU data}: synchronized phasor measurements at a high rate
    allow oscillations and events to be detected across the whole grid;
    spatiotemporal models (GNN + temporal networks) exploit the topology.
  \item \textbf{Distribution grids}: location is harder because of the
    branched radial topology and the lower level of instrumentation; ML
    combines data from smart meters, customer calls and reclosers.
\end{itemize}

Real fault data are scarce. For this reason, the models are often trained on
\textbf{electromagnetic transient simulations} (ATP/EMTP, PSCAD), which
creates the gap between simulation and reality (sim-to-real), one of the
main risks of these applications.

\subsection{Predictive maintenance of assets}

\paragraph{Power transformers.} They are among the most costly and critical
assets. \textbf{Dissolved gas analysis} (DGA) of the insulating oil is the
main diagnostic source: different types of internal fault (partial
discharges, arcing, overheating) generate characteristic proportions of
hydrogen, methane, ethylene, acetylene and other gases. Classical methods,
such as the Duval triangle, are graphical rules; ML allows DGA to be
combined with load history, temperature, age and tests to produce a health
index and a failure probability.

\paragraph{Wind turbines.} The SCADA systems of the turbines record, every
10 minutes, dozens of variables (bearing and gearbox temperatures, power,
wind, rotation). A widely used technique is the \textbf{normal behavior
model}:
\begin{enumerate}
  \item a model is trained to predict one variable (for example, bearing
    temperature) from the others, using only healthy periods;
  \item in operation, the residual $r_t = y_t - \hat{y}_t$ is monitored;
  \item a persistent deviation of the residual indicates degradation, often
    weeks before the failure.
\end{enumerate}
A ``persistent deviation'' needs a defined statistic, because the threshold
fixes the false alarm rate. A $3\sigma$ threshold applied to each sample,
with a Gaussian and independent residual, fires on 0.27\% of the samples:
0.39 false alarms per day per variable (144 samples of 10 minutes) and about
390 per day in a fleet of 100 turbines with 10 monitored variables.
Statistics that accumulate the deviation over time (CUSUM, EWMA) reduce this
rate for the same detection power.
High-frequency vibration sensors allow a finer diagnosis (bearings, gears)
with CNNs on spectra.

\paragraph{Photovoltaic plants.} Detection of underperforming strings by
comparison with neighboring strings and with the expected irradiance;
inverter failures; degradation and soiling (which also guides the cleaning
schedule).

\paragraph{Remaining useful life (RUL).} When there are enough failures in
the history, the time to failure can be estimated with survival models (Cox,
random survival forests, survival networks) or direct regression with
recurrent networks. The result feeds the optimization of the maintenance
plan, which weighs the cost of the intervention, the failure probability and
the consequence.

\subsection{Computer vision for inspection}

\begin{itemize}
  \item \textbf{Transmission and distribution lines}: detection of damaged
    insulators, missing dampers, corrosion and nests in images from drones
    and helicopters, with object detectors (YOLO family, DETR).
  \item \textbf{Vegetation}: LiDAR point clouds and satellite images to
    measure the distance between canopies and conductors and to prioritize
    pruning; contact with vegetation is one of the main causes of failures
    in distribution.
  \item \textbf{Thermography}: hot spots in connections, switches and solar
    panels.
  \item \textbf{Electroluminescence}: microcracks in photovoltaic cells.
\end{itemize}
The gain is one of scale: an inspection that would require crews covering
thousands of kilometers becomes an image screening problem, with humans
reviewing only the flagged cases.

\subsection{Data challenges}
\begin{itemize}
  \item \textbf{Extreme imbalance}: failures are rare. Overall accuracy is
    not informative; precision, recall, the PR curve and the cost of false
    positives (unnecessary visit) versus false negatives (unpredicted
    failure) are used instead. Precision depends on the prevalence $\pi$:
    \begin{equation}
      \text{precision} = \frac{\pi\,\mathrm{TPR}}
        {\pi\,\mathrm{TPR} + (1 - \pi)\,\mathrm{FPR}}.
    \end{equation}
    With $\pi = 1\%$, a recall (TPR) of 90\% and a false positive rate (FPR)
    of 5\%, the precision is 15\%: of every 100 alerts, 85 are false.
  \item \textbf{Poor labels}: work orders in free text, causes recorded
    inconsistently. Natural language processing (today with LLMs) helps to
    structure the maintenance history.
  \item \textbf{Scattered data}: SCADA, maintenance systems, oil laboratory
    and inspections in different systems, without a common key.
  \item \textbf{Asset heterogeneity}: different manufacturers, models and
    ages reduce the volume of comparable data. Transfer learning and
    semi-supervised models help.
\end{itemize}

\visaoproduto{Asset health platform (APM)}

\paragraph{Problem.}
Transmission companies, distribution utilities and generators maintain large
fleets of assets with maintenance plans based mainly on the calendar. This
produces unnecessary maintenance on healthy assets and unexpected failures
on degraded assets, with repair costs, lost revenue, continuity compensation
and safety risk.

\paragraph{Users and personas.}
\begin{itemize}
  \item \textbf{Maintenance engineer}: wants to know which assets need
    attention and why;
  \item \textbf{Asset manager}: wants to prioritize the maintenance and
    replacement budget based on risk;
  \item \textbf{Inspection crew}: wants prioritized routes and work orders.
\end{itemize}

\paragraph{Value proposition.}
A health and risk index per asset, updated continuously, with the
explanation of the signals behind it, integrated with the work order system.
The decision moves from ``maintain everything every $N$ months'' to
``maintain what is degrading, when it is degrading''.

\paragraph{MVP scope.}
Start with \textbf{one asset class} for which the data and the value are
established:
\begin{itemize}
  \item \textbf{Option A (wind generation)}: normal behavior models on
    turbine SCADA for gearbox, generator and bearings, with alerts and a
    dashboard per farm;
  \item \textbf{Option B (transmission/distribution)}: transformer health
    index from DGA, loading, age and test history.
\end{itemize}
In both cases: ranking of assets by risk, explanation of the alerts and
recording of the engineer's feedback (useful or false alert), which becomes
a label for the next training cycle.

\paragraph{Evolution (roadmap).}
\begin{itemize}
  \item \textbf{V2}: computer vision module for drone inspection (lines,
    substations, solar panels).
  \item \textbf{V3}: estimation of remaining useful life and optimization of
    the maintenance plan under budget and crew constraints.
  \item \textbf{V4}: automatic analysis of fault records for fault
    classification and location, integrated with restoration.
\end{itemize}

\paragraph{Reference architecture.}
\begin{itemize}
  \item \textbf{Ingestion} from SCADA/historians (such as PI), laboratory
    reports, maintenance management systems and images.
  \item \textbf{Unified asset register} (the key that links the data, in
    practice the largest task of the project).
  \item \textbf{Models} per asset class, with a common layer of alerts and
    explainability.
  \item \textbf{Workflow}: alert $\rightarrow$ engineer's analysis
    $\rightarrow$ work order $\rightarrow$ inspection result
    $\rightarrow$ label.
\end{itemize}

\paragraph{Success metrics.}
\begin{itemize}
  \item Technical: alert precision (proportion of confirmed alerts), mean
    lead time between alert and failure, failure recall.
  \item Operational: reduction of unscheduled failures, reduction of
    unnecessary maintenance, improvement of DEC/FEC and reduction of
    compensation.
  \item Financial: maintenance cost per asset and generation availability.
\end{itemize}

\paragraph{Business model.}
Subscription per monitored asset (per turbine, per transformer) plus a
deployment service. Outcome-based contracts (part of the value tied to the
reduction of failures) are possible when there is a defined baseline.

\paragraph{Risks.}
\begin{itemize}
  \item \textbf{Alert fatigue}: many false positives lead users to ignore
    the system. The precision of the alerts weighs more than their number.
  \item Quality and integration of legacy data.
  \item Competition from manufacturers (OEMs), which offer monitoring of
    their own equipment; the differentiation is to be multi-vendor.
\end{itemize}

% ===========================================================================
\section{Distribution and the end consumer}
% ===========================================================================

\subsection{Context}

Distribution is the segment in which the transformation is most visible:
MMGD grew after the legal framework for distributed generation (Law
14.300/2022), smart metering is advancing, and new services (time-of-use
tariffs, demand response, electric vehicles, a free market for smaller
consumers) change the relationship with the consumer. At the same time, old
problems, such as non-technical losses, continue to carry a large financial
weight for Brazilian distribution utilities.

\subsection{Non-technical losses}

\paragraph{The problem.} The losses of a distribution utility are the
difference between the energy injected into the grid and the energy billed.
Part of them are \emph{technical losses} (Joule effect in conductors and
transformers); the remainder are \emph{non-technical losses}: theft (illegal
connections), fraud (meter tampering), and metering and registration errors.
In the tariff reviews, ANEEL defines regulatory loss levels that can be
passed through to the tariff; whatever exceeds that level is a loss for the
utility. In some concession areas, non-technical losses are a relevant
fraction of the low-voltage market.

\paragraph{Formulation as classification.} For each consumer unit $u$, the
probability of irregularity
$p_u = P(\text{irregular} \mid \mathbf{x}_u)$ is estimated from attributes
such as:
\begin{itemize}
  \item \textbf{consumption profile}: abrupt drops, zero consumption, low
    variance, deviation from neighbors with the same profile;
  \item \textbf{history}: previous inspections, confirmed irregularities,
    meter replacements, complaints;
  \item \textbf{registration}: class, activity, declared load, meter type;
  \item \textbf{spatial context}: irregularity rate in the neighborhood and
    at the transformer;
  \item \textbf{energy balance}: when there is metering at the transformer,
    \begin{equation}
      E^{\text{NTL}}_{k} = E^{\text{transformer}}_{k} - \sum_{u \in k} E^{\text{metered}}_{u}
      - \hat{E}^{\text{technical}}_{k},
    \end{equation}
    indicates how much is lost under each transformer $k$ and helps to
    locate where to look. The balance is a small difference between large
    terms, and its uncertainty adds those of the terms:
    \begin{equation}
      \sigma^2\big(E^{\text{NTL}}_{k}\big) = \sigma^2_{\text{transformer}}
      + \sum_{u \in k} \sigma^2_{u} + \sigma^2_{\text{technical}}.
    \end{equation}
    The technical loss varies with the square of the current and depends on
    the shape of the load curve, not only on the energy of the period. The
    sum over $u \in k$ also assumes that the link between consumer and
    transformer is correct.
\end{itemize}
The output is a prioritized list of inspection targets.

\paragraph{Selection bias.} The labels come from inspections, and
inspections are not random: historically, they were carried out where there
was suspicion. The model therefore learns to reproduce the old selection
criteria, and not necessarily to identify fraud~\cite{ntlreview}. In formal
terms, training estimates
$P(\text{irregular} \mid \mathbf{x}_u, \text{inspected})$, which equals $p_u$
only if, given $\mathbf{x}_u$, the choice of whom to inspect does not depend
on the irregularity. Mitigations:
\begin{itemize}
  \item reserving part of the inspections for random or exploratory sampling
    (the exploration--exploitation balance);
  \item learning from positive and unlabeled examples (PU learning), treating
    the uninspected units as unlabeled and not as regular. The standard form
    of the method assumes that the labeled positives are a random sample of
    the positives, which selection by suspicion violates; the applicable
    variant models the probability of inspection as a function of
    $\mathbf{x}_u$ and weights the examples by it;
  \item unsupervised anomaly detection models as a complement.
\end{itemize}

\paragraph{Metrics.} Because inspection capacity is limited, the relevant
metrics are the \textbf{hit rate at the top of the list} (precision at $k$)
and the \textbf{energy recovered per inspection}, and not global metrics
such as AUC. The two metrics call for different orderings. Precision at $k$
is maximized by ordering by $p_u$; the energy recovered is maximized by
ordering by the expected value $p_u\,\hat{E}_u$, where $\hat{E}_u$ is the
estimated recoverable energy of the unit. The second ordering requires a
calibrated $p_u$, because the product uses the value of the probability and
not only its rank.

\subsection{Load disaggregation (NILM)}

Non-intrusive load monitoring estimates the consumption of each appliance
from a single aggregate measurement~\cite{hart}:
\begin{equation}
  P(t) = \sum_{i=1}^{N} p_i(t) + \varepsilon(t).
\end{equation}
It is an ill-posed source separation problem: many combinations of
appliances explain the same aggregate.

\paragraph{Methods.}
\begin{itemize}
  \item \textbf{Classical}: event detection (power steps) and factorial
    hidden Markov models.
  \item \textbf{Deep learning}: networks that map a window of the aggregate
    to the consumption of one appliance, either the whole window (seq2seq)
    or the central point (seq2point)~\cite{kelly,seq2point}. They are trained
    on public datasets such as REDD and UK-DALE.
\end{itemize}

\paragraph{Practical limitation.} Quality depends strongly on the sampling
rate. With meters that record every 15 minutes or 1 hour, only loads with
high power and prolonged use can be identified (air conditioning, heater,
electric vehicle charger). A high and short load is diluted in the interval
mean: a 5\,kW electric shower switched on for 5 minutes appears as 1.67\,kW
in a 15-minute window. Fine disaggregation requires data at a resolution of
seconds or finer, obtained with additional devices.

\subsection{Distribution grid with MMGD and electric vehicles}

\begin{itemize}
  \item \textbf{Low-voltage forecasting}: series per transformer or per
    consumer are noisy and hard to forecast individually. Probabilistic peak
    forecasting, including with foundation models, is an approach under
    evaluation~\cite{lvpeak}.
  \item \textbf{Hosting capacity}: how much MMGD or how many chargers a
    feeder supports without violating voltage and loading limits. ML
    surrogates accelerate studies that would require thousands of
    simulations.
  \item \textbf{Topology and phase identification}: records of which
    consumer is on which transformer and phase are frequently wrong; the
    correlation between voltage series from smart meters allows them to be
    corrected.
  \item \textbf{Detection of unregistered MMGD} from the consumption
    profile.
\end{itemize}

\subsection{Demand response and efficiency}

\begin{itemize}
  \item \textbf{Time-of-use tariffs}: with the Tarifa Branca (the Brazilian
    time-of-use tariff for low-voltage consumers) and, in the future, more
    granular tariffs, the consumer needs to understand when to shift
    consumption. Forecasting and recommendation models personalize this
    guidance.
  \item \textbf{Control of flexible loads}: air conditioning, water heating,
    vehicle charging and batteries can be scheduled by model predictive
    control (MPC) or RL to minimize cost while respecting comfort.
  \item \textbf{Commercial buildings and industry}: consumption forecasting
    and anomaly detection identify waste (equipment switched on outside
    working hours, degradation of HVAC systems). The DeepMind case, with a
    reduction of up to 40\% in the cooling energy of data centers,
    illustrates the potential~\cite{deepmindcooling}.
  \item \textbf{Customer segmentation}: clustering of load curves for tariff
    design, efficiency campaigns and offers in the free market.
\end{itemize}

\subsection{Privacy and regulation}

High-resolution consumption data reveal habits (schedules, presence,
appliances). In Brazil, they are personal data under the LGPD (Law
13.709/2018). This requires a legal basis for processing, minimization,
transparency and security. Techniques such as aggregation, federated
learning and differential privacy reduce the exposure. The rollout of smart
metering in Brazil is also more gradual than in Europe, which limits the
availability of high-resolution data at large scale.

\visaoproduto{Loss and grid intelligence for distribution utilities}

\paragraph{Problem.}
Distribution utilities lose revenue to non-technical losses above the
regulatory level and carry out inspections with a low hit rate. At the same
time, they need to handle an increasingly complex low-voltage grid (MMGD,
electric vehicles) with inaccurate records and little observability.

\paragraph{Users and personas.}
\begin{itemize}
  \item \textbf{Energy recovery manager}: wants to maximize the energy
    recovered with the available inspection budget;
  \item \textbf{Field crew}: wants prioritized routes and targets, with
    low-effort recording of the result;
  \item \textbf{Distribution planning engineer}: wants to know where the
    grid will saturate with MMGD and new loads.
\end{itemize}

\paragraph{Value proposition.}
More energy recovered per inspection, with a prioritized and explained list
of targets that improves at each cycle with the results of the inspections
themselves. Afterwards, the same database supports the visibility of the
low-voltage grid.

\paragraph{MVP scope.}
\begin{enumerate}
  \item Model for prioritizing inspection targets with billing,
    registration and inspection history data.
  \item Explanation per target (why this customer was indicated).
  \item Closed loop: the inspection result returns as a label; part of the
    inspections is reserved for exploration, which corrects the selection
    bias.
  \item Results dashboard: hit rate and energy recovered compared with the
    utility's current method (A/B test by region or by batch).
\end{enumerate}

\paragraph{Evolution (roadmap).}
\begin{itemize}
  \item \textbf{V2}: energy balance per transformer where there is metering,
    combining location (where the loss occurs) with prioritization (whom to
    investigate).
  \item \textbf{V3}: correction of topology and phase records, detection of
    unregistered MMGD, forecasting of transformer loading.
  \item \textbf{V4}: hosting capacity studies and a consumer-facing product
    (disaggregation, tariff and efficiency recommendations).
\end{itemize}

\paragraph{Reference architecture.}
\begin{itemize}
  \item \textbf{Ingestion} from the commercial system (billing,
    registration), the work order management system, smart metering where
    available, and georeferenced grid data.
  \item \textbf{Prioritization engine} with supervised models (gradient
    boosting is the standard) and an unsupervised anomaly component.
  \item \textbf{Field application} to record inspection results in a
    structured form.
  \item \textbf{Data governance} in accordance with the LGPD (access
    control, anonymization in aggregate analyses, record of processing).
\end{itemize}

\paragraph{Success metrics.}
\begin{itemize}
  \item Technical: precision at the top of the list; calibration of the
    probabilities.
  \item Business: energy recovered (MWh and R\$) per inspection; reduction
    of the non-technical loss index; cost per productive inspection.
  \item Adoption: percentage of inspections generated by the system;
    percentage of results recorded.
\end{itemize}

\paragraph{Business model.}
Subscription plus a variable component tied to the energy recovered (success
fee). The gain is measurable in reais and appears early, which helps the
sale. R\&D contracts regulated by ANEEL can also fund the grid modules that
are still experimental.

\paragraph{Risks.}
\begin{itemize}
  \item Label quality and selection bias; without an exploratory cycle the
    model becomes self-referential.
  \item Social and legal sensitivity: the indication of an irregularity
    needs a fair procedure; the model prioritizes inspections and does not
    constitute evidence.
  \item Geographic or socioeconomic bias, which needs to be monitored to
    avoid an unjustified concentration of inspections on certain groups.
  \item Compliance with the LGPD and security of consumer data.
\end{itemize}

% ===========================================================================
\section{Comparative summary}
% ===========================================================================

\begin{table}[h]
\small
\caption{Comparison of the four areas.}
\begin{tabelalarga}
\begin{tabularx}{\linewidth}{l L L L l}
\toprule
\textbf{Area} & \textbf{Core techniques} & \textbf{Product} &
  \textbf{Main customer} & \textbf{Maturity} \\
\midrule
Forecasting & GBM, LSTM, Transformers, foundation models &
  Forecasting platform & generators, trading companies, distribution
  utilities & high \\
Grid operation & GNN, PINN, safe RL & Operation copilot &
  transmission companies, distribution utilities, renewable operators &
  medium/low \\
Faults and maintenance & CNN, normal behavior models, survival &
  Asset health platform & generators, transmission companies, distribution
  utilities & high \\
Distribution and consumer & GBM, PU learning, NILM, MPC &
  Loss and grid intelligence & distribution utilities & high (losses),
  medium (grid) \\
\bottomrule
\end{tabularx}
\end{tabelalarga}
\end{table}

\paragraph{Patterns common to the four areas.}
\begin{enumerate}
  \item \textbf{Data before models}: in all areas, the largest deployment
    effort is integrating and cleaning data from legacy systems, not
    training networks.
  \item \textbf{Physics as an ally}: models that respect physical laws
    (PINNs, verification by solver, normalizations such as the clear-sky
    index) tend to generalize better and to be more readily accepted by
    engineers.
  \item \textbf{Human in control}: the products that are viable in the short
    term support decisions instead of making them alone, and they learn
    from user feedback.
  \item \textbf{Shadow mode as a sales strategy}: running in parallel and
    showing the gain with the customer's own data reduces the perceived
    risk.
  \item \textbf{Suggested order of entry}: forecasting and non-technical
    losses have early and measurable value and serve as the entry point;
    predictive maintenance comes next; the operation copilot is the bet
    with the longest horizon and the highest barrier to entry.
\end{enumerate}

% ===========================================================================
% References
% ===========================================================================
\begin{thebibliography}{99}
\small

\bibitem{energyinformatics2025}
Machine learning applications in energy systems: current trends, challenges,
and research directions. \textit{Energy Informatics}, 2025.
\url{https://energyinformatics.springeropen.com/articles/10.1186/s42162-025-00524-6}

\bibitem{trustworthy2024}
Trustworthy artificial intelligence in the energy sector: landscape analysis
and evaluation framework. arXiv:2412.07782, 2024.

\bibitem{deepmindwind}
DeepMind. Machine learning can boost the value of wind energy, 2019. Case
summarized at
\url{https://bestpractice.ai/ai-use-cases/case-studies/energy-utilities/deepmind-increases-value-of-wind-power-by-20-by-predicting-supply-36-hours-in-advance}

\bibitem{deepmindcooling}
DeepMind. DeepMind AI reduces Google data centre cooling bill by 40\%, 2016.
\url{https://deepmind.google/blog/deepmind-ai-reduces-google-data-centre-cooling-bill-by-40/}

\bibitem{onssolar}
ONS. Avaliação e diagnóstico do sistema de previsão da irradiância solar e
geração. Workshop presentation.
\url{https://www.ons.org.br/AcervoDigitalDocumentosEPublicacoes/Apresentacao_Workshop.pdf}

\bibitem{onseolico}
Desenvolvimento e implantação no ONS de um modelo de previsão de geração
eólica. ABEEólica.
\url{https://abeeolica.org.br/wp-content/uploads/2017/07/Art.pdf}

\bibitem{pvlstm}
Previsão horária de geração fotovoltaica utilizando deep learning.
\textit{Observatorio de la Economía Latinoamericana}.
\url{https://ojs.observatoriolatinoamericano.com/ojs/index.php/olel/article/view/11654}

\bibitem{tft}
B.~Lim, S.~Ö.~Arık, N.~Loeff, T.~Pfister. Temporal Fusion Transformers for
interpretable multi-horizon time series forecasting. \textit{International
Journal of Forecasting}, 2021.

\bibitem{chronos}
A.~F.~Ansari et al. Chronos: learning the language of time series.
\textit{Transactions on Machine Learning Research}, 2024.

\bibitem{timesfm}
A.~Das, W.~Kong, R.~Sen, Y.~Zhou. A decoder-only foundation model for
time-series forecasting. \textit{ICML}, 2024.

\bibitem{fmreview}
Foundation models for clean energy forecasting: a comprehensive review.
arXiv:2507.23147, 2025.

\bibitem{loadbenchmark}
A benchmark for electrical load forecasting across grid levels: time-series
transformers outperform established methods. arXiv:2607.15705, 2026.

\bibitem{tsfmconsumer}
Time series foundation models for energy load forecasting on consumer
hardware: a multi-dimensional zero-shot benchmark. arXiv:2602.10848, 2026.

\bibitem{tsfmxai}
Explainable load forecasting with covariate-informed time series foundation
models. arXiv:2604.28149, 2026.

\bibitem{lvpeak}
Probabilistic low-voltage peak load forecasting with time series foundation
models evaluated on application-oriented metrics. arXiv:2607.01966, 2026.

\bibitem{gridfm}
GridFM: a physics-informed foundation model for multi-task energy forecasting
using real-time NYISO data. \textit{Energies}, 19(2):357, 2026.
\url{https://doi.org/10.3390/en19020357}

\bibitem{gnnpf}
Graph neural networks for efficient AC power flow prediction in power grids.
arXiv:2502.05702, 2025.

\bibitem{gnnopfinit}
Initial estimate of AC optimal power flow with graph neural networks.
\textit{Electric Power Systems Research}, 2024.

\bibitem{gnnreconf}
Graph neural network-accelerated network-reconfigured optimal power flow.
arXiv:2410.17460, 2024.

\bibitem{dc3}
P.~L.~Donti, D.~Rolnick, J.~Z.~Kolter. DC3: a learning method for optimization
with hard constraints. \textit{ICLR}, 2021.

\bibitem{l2rpn}
A.~Marot et al. Learning to run a power network challenge for training
topology controllers. \textit{Electric Power Systems Research}, 2020.

\bibitem{graphrlsurvey}
Graph reinforcement learning for power grids: a comprehensive survey.
\textit{Energy and AI}, 2025.

\bibitem{ppognn}
Proximal policy optimization with graph neural networks for optimal power
flow. arXiv:2212.12470, 2022.

\bibitem{rlpinn}
Optimizing energy management of smart grid using reinforcement learning aided
by surrogate models built using physics-informed neural networks.
arXiv:2510.17380, 2025.

\bibitem{saferl}
A review of safe reinforcement learning methods for modern power systems.
arXiv:2407.00304, 2024.

\bibitem{gridgeneralization}
Towards systematic generalization for power grid optimization problems.
arXiv:2605.02026, 2026.

\bibitem{pinn}
M.~Raissi, P.~Perdikaris, G.~E.~Karniadakis. Physics-informed neural
networks. \textit{Journal of Computational Physics}, 378:686--707, 2019.

\bibitem{faultsurvey}
Survey on methods for detection, classification and location of faults in
power systems using artificial intelligence. arXiv:2507.10011, 2025.

\bibitem{ntlreview}
P.~Glauner et al. The challenge of non-technical loss detection using
artificial intelligence: a survey. \textit{International Journal of
Computational Intelligence Systems}, 2017.

\bibitem{hart}
G.~W.~Hart. Nonintrusive appliance load monitoring. \textit{Proceedings of
the IEEE}, 80(12):1870--1891, 1992.

\bibitem{kelly}
J.~Kelly, W.~Knottenbelt. Neural NILM: deep neural networks applied to energy
disaggregation. \textit{ACM BuildSys}, 2015.

\bibitem{seq2point}
C.~Zhang, M.~Zhong, Z.~Wang, N.~Goddard, C.~Sutton. Sequence-to-point
learning with neural networks for non-intrusive load monitoring.
\textit{AAAI}, 2018.

\end{thebibliography}

\end{document}
